% !TEX program = xelatex
% =====================================================================
%  BÁO CÁO ĐỒ ÁN TRÍ TUỆ NHÂN TẠO, HỌC PHẦN AIP202
%  Khoa Công nghệ Thông tin, Trường Đại học Kiến trúc Đà Nẵng
%
%  Build:  latexmk -xelatex do-an.tex
%  Overleaf: Menu > Compiler > XeLaTeX, rồi bấm Recompile.
%
%  Khung tài liệu chuyển thể từ mẫu đồ án của Trường Đại học Nha Trang
%  (github.com/nd-hung/thesis-template, tác giả TS. Nguyễn Đình Hưng).
%  Toàn bộ nhận diện, phông chữ, bảng màu, trang bìa và cấu trúc chương
%  được thay theo chuẩn của DAU và theo đề cương học phần AIP202.
% =====================================================================
% \documentclass phải nằm ở file này, không được chuyển vào settings.tex,
% vì Overleaf dò file chính bằng cách tìm file có chứa \documentclass.
\documentclass[12pt, a4paper, oneside]{book}
% !TEX root = ../do-an.tex
% =====================================================================
%  styles/settings.tex
%  Các thiết lập chung, nạp từ file chính do-an.tex.
%  SINH VIÊN KHÔNG SỬA FILE NÀY.
%
%  Lưu ý: lệnh khai báo lớp tài liệu nằm ở file do-an.tex chứ không nằm
%  ở đây. Lý do: Overleaf dò file chính bằng cách tìm chuỗi khai báo lớp
%  tài liệu trong các file. Nếu chuỗi đó xuất hiện ở file này, Overleaf
%  sẽ biên dịch nhầm file này và báo lỗi "no legal end found", vì file
%  này chỉ chứa thiết lập chứ không chứa phần thân tài liệu.
% =====================================================================

% Toàn bộ nhận diện DAU: phông, màu, bìa, hộp trình bày, header
\usepackage{styles/dau}
\dauLayout

% Giãn dòng và khoảng cách đoạn theo quy cách đồ án
\linespread{1.25}
\usepackage[skip=8pt plus 1pt, indent=15pt]{parskip}

% Ghi chú: mẫu gốc dùng gói atbegshi để vứt trang đầu, vì trang bìa của
% mẫu gốc dựng bằng minipage nên sinh ra một trang trắng. Trang bìa ở
% đây dựng bằng môi trường titlepage nên không có trang trắng, và nếu
% giữ thủ thuật đó thì chính trang bìa sẽ bị vứt mất.


% =====================================================================
%  TỐI ƯU THỜI GIAN BIÊN DỊCH (TRÁNH LỖI TIMEOUT)
%  Khi hoàn thiện báo cáo, để các dòng \includeonly được comment để
%  biên dịch đầy đủ toàn bộ các chương.
% =====================================================================
%\includeonly{chuong/01-mo-dau}
%\includeonly{chuong/02-co-so-ly-thuyet}
%\includeonly{chuong/03-phan-tich-thiet-ke}
%\includeonly{chuong/04-cai-dat}
%\includeonly{chuong/05-ket-qua-danh-gia}
%\includeonly{chuong/06-ket-luan}

% =====================================================================
%  CẤU HÌNH TIÊU ĐỀ TIẾNG VIỆT CHO CÁC DANH MỤC
% =====================================================================
\renewcommand{\contentsname}{Mục lục}
\renewcommand{\listfigurename}{Danh sách hình vẽ}
\renewcommand{\listtablename}{Danh sách bảng}
\providecommand{\listalgorithmname}{Danh mục giải thuật}
\providecommand{\listalgorithmcfname}{Danh mục giải thuật}
\renewcommand{\listalgorithmname}{Danh mục giải thuật}
\renewcommand{\listalgorithmcfname}{Danh mục giải thuật}

% =====================================================================
%  PHẦN EM ĐIỀN: thông tin đồ án
% =====================================================================
% Loại báo cáo
\crname{BÁO CÁO ĐỒ ÁN TRÍ TUỆ NHÂN TẠO}
% Học phần
\cmonhoc{Học phần AIP202 \textperiodcentered{} 2 tín chỉ}
% Tên đề tài. Dùng \\ để ngắt dòng cho cân đối trên bìa.
\ctname{XÂY DỰNG TRỢ LÝ AI CÁ NHÂN HÓA\\HỖ TRỢ NGƯỜI KHIẾM THỊ}

% Họ tên sinh viên và mã số sinh viên
\cstuname{Phạm Phước Trung}
\cmssv{2251220208}

\cstunameTwo{Phạm Hoàng Nam}
\cmssvTwo{2151220151}
% Lớp
\clopsv{CNTT 23CT2}
% Ngành
\csCouncil{Công nghệ thông tin}
% Giảng viên hướng dẫn
\csSupervise{TS. Đỗ Phúc Hảo}
% Năm học
\cttime{2026--2027}
% Khoa quản lý
\csdeptname{KHOA CÔNG NGHỆ THÔNG TIN}
% Nơi thực hiện
\csplace{Đà Nẵng}

\begin{document}

% Bìa trang trọng có khung kép. Muốn dùng bìa hiện đại theo bộ giáo
% trình của khoa thì thay dòng dưới bằng \coverpagehiendai
\coverpage

% ---------------------------------------------------------------------
%  PHẦN ĐẦU: đánh số trang bằng số La Mã thường
% ---------------------------------------------------------------------
\frontmatter
\dauPageStyle

% ---------------------------------------------------------------------
%  LỜI CAM ĐOAN
% ---------------------------------------------------------------------
\begin{declaration}

Chúng tôi xin cam đoan đồ án này là công trình do chính chúng tôi thực hiện dưới sự
hướng dẫn của \dauGVHD. Các nội dung nghiên cứu, số liệu và kết quả trình
bày trong đồ án là trung thực và chưa từng được công bố trong bất kỳ công
trình nào khác.

Mọi số liệu, hình ảnh, mã nguồn và kết quả của tác giả khác được sử dụng
trong đồ án đều được trích dẫn nguồn gốc rõ ràng trong phần tài liệu tham
khảo. Toàn bộ phần mềm và thư viện được dùng đều có giấy phép cho phép sử
dụng. Những phần có sử dụng công cụ trí tuệ nhân tạo hỗ trợ đã được nêu rõ
trong báo cáo.

Nếu phát hiện có bất kỳ sự gian lận nào, chúng tôi xin chịu hoàn toàn trách nhiệm
trước Khoa và Nhà trường.

\vspace{1.4em}
\hfill
\begin{tabular}{c}
\textbf{Sinh viên thực hiện} \\[3.5em]
\textbf{Phạm Phước Trung} \quad \textbf{Phạm Hoàng Nam}
\end{tabular}

\end{declaration}
% ---------------------------------------------------------------------
%  LỜI CẢM ƠN
% ---------------------------------------------------------------------
\begin{acknowledgments}

Chúng em xin gửi lời cảm ơn chân thành đến \dauGVHD, người đã tận tình định hướng
đề tài, góp ý và theo sát tiến độ của chúng em trong suốt quá trình thực hiện đồ
án. Chúng em cũng xin cảm ơn quý thầy cô Khoa Công nghệ Thông tin, Trường Đại học
Kiến trúc Đà Nẵng đã trang bị cho chúng em nền tảng kiến thức để hoàn thành công
việc này. Cuối cùng, chúng em xin cảm ơn gia đình và bạn bè đã luôn động viên chúng em
trong thời gian học tập.
\vspace{1.4em}
\hfill
\begin{tabular}{c}
\textbf{Sinh viên thực hiện} \\[3.5em]
\textbf{Phạm Phước Trung} \quad \textbf{Phạm Hoàng Nam}
\end{tabular}

\end{acknowledgments}
% ---------------------------------------------------------------------
%  TÓM TẮT
% ---------------------------------------------------------------------
\begin{abstract}

Đồ án xây dựng VisionAI Pro, một nguyên mẫu trợ lý thị giác AI hỗ trợ người khiếm thị tiếp cận thông tin trong ảnh bằng tiếng Việt. Ứng dụng Streamlit cho phép tải ảnh hoặc chụp ảnh để Gemini mô tả bối cảnh, vật thể, nguy cơ, chữ và tiền Việt; chế độ OCR riêng sử dụng Tesseract cục bộ. Người dùng có thể hỏi đáp theo ảnh, tìm vật trong ảnh, dùng lệnh giọng nói để điều hướng hoặc đọc lại kết quả, và nghe phản hồi bằng Edge TTS hoặc Windows SAPI5. Giao diện bổ sung robot Vision tương tác, chế độ trợ năng, lịch sử cục bộ tối đa 100 bản ghi, chức năng lọc/yêu thích, thống kê, xuất báo cáo TXT và kiểm tra kết nối Gemini. STT sử dụng SpeechRecognition/Google Web Speech và cần Internet; các tuỳ chọn cá nhân hoá chỉ duy trì trong phiên làm việc. Trên bộ source được đóng gói, 15/15 kiểm thử unit đạt và endpoint health của Streamlit trả về \texttt{ok}. Chưa kiểm thử gọi Gemini thật, chưa đo độ chính xác/độ trễ và chưa đánh giá với người dùng; các kết quả phần mềm này không chứng minh chất lượng hoặc độ an toàn của AI. Hệ thống chỉ phân tích ảnh tĩnh, không thay thế thiết bị hỗ trợ di chuyển đã được kiểm chứng.

\vspace{0.8em}
\noindent\textbf{Từ khoá:} VisionAI Pro, Gemini Vision, OCR, trợ lý thị giác, khả năng tiếp cận.

\end{abstract}


% --- MỤC LỤC ---
\cleardoublepage
\phantomsection
\addcontentsline{toc}{chapter}{Mục lục}
\tableofcontents

% --- DANH SÁCH HÌNH VẼ ---
\cleardoublepage
\phantomsection
\addcontentsline{toc}{chapter}{Danh sách hình vẽ}
\listoffigures

% --- DANH SÁCH BẢNG ---
\cleardoublepage
\phantomsection
\addcontentsline{toc}{chapter}{Danh sách bảng}
\listoftables

% --- DANH MỤC GIẢI THUẬT ---
\cleardoublepage
\phantomsection
\addcontentsline{toc}{chapter}{Danh mục giải thuật}
\listofalgorithms

% ---------------------------------------------------------------------
%  THÂN BÁO CÁO: sáu chương theo quy cách của học phần
% ---------------------------------------------------------------------
\mainmatter
\dauPageStyle

% ---------------------------------------------------------------------
%  CHƯƠNG 1 — MỞ ĐẦU
% ---------------------------------------------------------------------
\chapter{Mở đầu}
\label{ch:modau}

\begin{muctieu}
Chương mở đầu trình bày bối cảnh, lý do lựa chọn đề tài, mục tiêu, phạm vi
và đối tượng thực hiện của đồ án. Nội dung chương nhằm làm rõ nhu cầu xây
dựng một trợ lý AI có khả năng cá nhân hoá để hỗ trợ người khiếm thị trong
quá trình tiếp cận thông tin, nhận diện không gian/vật thể và tương tác với hệ thống.
\end{muctieu}

\section{Lý do chọn đề tài}

Trong quá trình phát triển của trí tuệ nhân tạo, các trợ lý AI ngày càng
được ứng dụng rộng rãi để hỗ trợ con người trong việc tìm kiếm thông tin,
giao tiếp và thực hiện các tác vụ hằng ngày. Tuy nhiên, phần lớn các trợ lý
AI hiện nay được thiết kế theo hướng phục vụ người dùng nói chung, chủ yếu dựa
trên giao diện trực quan và hình ảnh, chưa tập trung đầy đủ vào nhu cầu tiếp
cận thông tin và tương tác bằng âm thanh/giọng nói của người khiếm thị.

Người khiếm thị gặp rất nhiều khó khăn trong các tình huống giao tiếp và sinh
hoạt hằng ngày vốn phụ thuộc vào thị giác, chẳng hạn như đọc tài liệu văn bản,
nhận diện các vật thể xung quanh, quan sát biển hiệu hay di chuyển trong môi
trường lạ. Vì vậy, một hệ thống hỗ trợ dựa trên trí tuệ nhân tạo có khả năng
chuyển đổi thông tin thị giác thành âm thanh/lời nói hoặc phản hồi thông minh
qua giọng nói có thể giúp giảm bớt những trở ngại trong quá trình tiếp cận công nghệ
và hòa nhập cộng đồng~\cite{WHO2019WorldVision,WHO2022AssistiveTech}.

Bên cạnh khả năng cung cấp và hỗ trợ xử lý thông tin, khả năng điều chỉnh cách
tương tác cũng có ý nghĩa quan trọng. Mỗi người dùng khiếm thị có thể có nhu cầu
khác nhau về cỡ chữ, thao tác, giọng đọc và mức độ chi tiết. Trong phạm vi
nguyên mẫu này, các lựa chọn được áp dụng trong phiên làm việc; hệ thống chưa
lưu hồ sơ cá nhân lâu dài.

Xuất phát từ những vấn đề trên, đồ án lựa chọn đề tài ``Xây dựng trợ lý AI
cá nhân hoá hỗ trợ người khiếm thị'' nhằm nghiên cứu khả năng ứng dụng
trí tuệ nhân tạo vào một bài toán có tính thực tiễn cao, hướng đến việc gia tăng
khả năng tiếp cận công nghệ và nâng cao chất lượng cuộc sống cho người dùng có
nhu cầu đặc thù.

\section{Mục tiêu của đồ án}

Đồ án hướng tới các mục tiêu sau:
\begin{itemize}[leftmargin=1.6em, itemsep=2pt]
  \item Phân tích nhu cầu và các khó khăn trong quá trình tiếp cận thông tin,
        quan sát và tương tác với trợ lý AI của người khiếm thị.
        
  \item Xây dựng một trợ lý AI có khả năng tiếp nhận yêu cầu bằng giọng nói
        hoặc dữ liệu đầu vào (văn bản, hình ảnh), xử lý nội dung và đưa ra
        phản hồi bằng âm thanh/giọng nói phù hợp.
        
  \item Bổ sung tuỳ chọn tương tác trong phiên như chế độ trợ năng, giọng đọc,
        mức chi tiết và ngữ cảnh hội thoại; không đặt mục tiêu xây dựng hồ sơ
        người dùng lâu dài trong phiên bản này.
        
  \item Thiết kế giao diện dễ tiếp cận với chế độ chữ/nút lớn, focus bàn phím,
        robot hướng dẫn và khả năng nghe kết quả bằng giọng đọc.
        
  \item Kiểm thử các chức năng phần mềm đã xây dựng; xác định dữ liệu và phương pháp
        cần có để đánh giá riêng độ chính xác, độ trễ và mức độ đáp ứng bài toán.
\end{itemize}

\section{Phạm vi và đối tượng}

Đồ án tập trung xây dựng một nguyên mẫu trợ lý thị giác AI hướng đến hỗ trợ
người khiếm thị tiếp nhận thông tin trong ảnh, đọc văn bản và tương tác với
trợ lý bằng tiếng Việt. Đối tượng hướng đến là người mù hoặc người có thị lực
kém cần công cụ hỗ trợ đọc và nhận biết thông tin môi trường.

Phạm vi của đồ án tập trung vào nguyên mẫu Streamlit gồm phân tích ảnh Gemini,
OCR cục bộ bằng Tesseract, hỏi đáp theo ảnh, lệnh thoại, đọc phản hồi và các
tuỳ chọn trợ năng. STT sử dụng Google Web Speech qua Internet; TTS ưu tiên
Edge TTS và có Windows SAPI5 dự phòng. Các tuỳ chọn cá nhân hoá chỉ tồn tại
trong phiên làm việc, không gồm tài khoản hoặc hồ sơ lâu dài.

Đồ án không hướng tới việc thay thế hoàn toàn các công cụ hỗ trợ di chuyển chuyên dụng
(như gậy trắng, chó dẫn đường) hay các thiết bị y tế chuyên khoa, không giải quyết
toàn bộ các vấn đề y tế liên quan đến thị giác và chưa đặt mục tiêu triển khai hệ thống
ở quy mô thương mại/sản xuất lớn. Các chức năng và phạm vi hỗ trợ cụ thể được giới hạn
theo thời gian và nguồn lực thực hiện đồ án.

\begin{luuy}
Phạm vi của đồ án được giới hạn ở việc xây dựng và đánh giá một nguyên mẫu
trợ lý AI cá nhân hoá hỗ trợ người khiếm thị. Việc giới hạn phạm vi giúp
tập trung vào các chức năng cốt lõi, bảo đảm hệ thống có thể được triển khai,
kiểm thử và đánh giá trong thời gian thực hiện đồ án.
\end{luuy}

\section{Bố cục báo cáo}

Phần còn lại của báo cáo được tổ chức như sau. Chương~\ref{ch:coso} trình
bày cơ sở lý thuyết và các công nghệ được sử dụng trong quá trình xây dựng
trợ lý AI. Chương~\ref{ch:thietke} phân tích bài toán, yêu cầu hệ thống và
thiết kế giải pháp. Chương~\ref{ch:caidat} mô tả quá trình cài đặt và xây
dựng các chức năng của hệ thống. Chương~\ref{ch:ketqua} trình bày kết quả
thực nghiệm và đánh giá hệ thống. Chương~\ref{ch:ketluan} tổng kết những
kết quả đạt được, các hạn chế và hướng phát triển của đề tài.
% ---------------------------------------------------------------------
%  CHƯƠNG 2 — CƠ SỞ LÝ THUYẾT VÀ CÔNG NGHỆ
% ---------------------------------------------------------------------
\chapter{Cơ sở lý thuyết và công nghệ}
\label{ch:coso}

\begin{muctieu}
Chương này trình bày các cơ sở lý thuyết và công nghệ thực tế được sử dụng
trong hệ thống trợ lý thị giác AI. Nội dung tập trung vào AI đa phương thức,
thị giác máy tính, OCR, nhận dạng và tổng hợp giọng nói, giao diện Streamlit
và cơ chế lưu lịch sử/cấu hình phiên làm việc.
\end{muctieu}

\section{Kiến thức nền tảng}

\subsection{Trí tuệ nhân tạo và AI đa phương thức}

Trí tuệ nhân tạo \en{(Artificial Intelligence -- AI)} là lĩnh vực nghiên cứu
các phương pháp giúp máy tính thực hiện những nhiệm vụ thường cần khả năng
nhận thức và xử lý thông tin của con người. Trong đồ án, AI được sử dụng để
phân tích hình ảnh, nhận biết vật thể, bối cảnh, nguy cơ, chữ và mệnh giá tiền,
sau đó tạo phản hồi bằng tiếng Việt.

Hệ thống có tính chất đa phương thức vì tiếp nhận hình ảnh, văn bản và bản ghi
âm thanh đã được chuyển thành văn bản, đồng thời tạo ra văn bản và âm thanh phản
hồi. Thành phần Vision sử dụng Gemini thông qua Google GenAI SDK để phân tích
hình ảnh; API hỗ trợ đầu ra JSON có cấu trúc
~\cite{GoogleGeminiVision2026,GoogleGeminiStructuredOutput2026,GoogleGenAIPython2026}.
Trong source, Gemini được yêu cầu trả JSON qua MIME type; chương trình chuẩn hoá
các trường nhận được trước khi hiển thị. Thành phần giọng nói đảm nhiệm STT và TTS.

\subsection{Thị giác máy tính và Vision Language Model}

Thị giác máy tính \en{(Computer Vision)} cho phép máy tính trích xuất thông tin
có ý nghĩa từ hình ảnh. Trong hệ thống, ảnh được tải lên hoặc chụp từ camera,
sau đó được chuyển thành dữ liệu phù hợp để gửi đến Gemini Vision
~\cite{PillowImageModule2026}.

Mô hình Vision được yêu cầu trả về các trường có cấu trúc như mô tả, độ tin
cậy, bối cảnh, vật thể, nguy cơ, hướng dẫn di chuyển, văn bản nhận diện và
mệnh giá tiền. Cách trả kết quả dạng JSON giúp chương trình Python có thể xử
lý và hiển thị từng thành phần một cách ổn định.

Ở mức nền tảng, Transformer và Vision Transformer là các hướng nghiên cứu áp
dụng attention cho dữ liệu chuỗi và ảnh; CLIP là ví dụ về học liên kết ảnh với
ngôn ngữ tự nhiên. Đây là tài liệu tham khảo lý thuyết, không phải các mô hình
được gọi trực tiếp trong source của đồ án
~\cite{Vaswani2017Attention,Dosovitskiy2021ViT,Radford2021CLIP}.

Trong phiên bản hiện tại, việc chuẩn hoá ảnh trước khi gửi giúp giới hạn kích
thước dữ liệu truyền; kết quả JSON giúp giao diện truy cập các trường theo
tên thay vì phân tích văn bản tự do. Đây là lựa chọn về luồng dữ liệu, không
phải cam kết về thời gian phản hồi cố định.

\subsection{Nhận dạng ký tự quang học (OCR)}

OCR \en{(Optical Character Recognition)} là kỹ thuật chuyển nội dung chữ
trong hình ảnh thành văn bản. Đồ án sử dụng Tesseract OCR chạy cục bộ cho
chế độ ``Đọc chữ trong ảnh (OCR)''. Đây là một nhánh xử lý riêng, không giả
lập OCR bằng Gemini.

Trước khi nhận dạng, ảnh được chuyển sang RGB, phóng to khi cần, tăng tương
phản, tăng độ sắc nét và tạo thêm biến thể nhị phân. Tesseract được chạy với
các chế độ phân đoạn phù hợp; dữ liệu tiếng Việt \texttt{vie} và tiếng Anh \texttt{eng} được
ưu tiên khi có sẵn. Kết quả được lọc theo độ tin cậy và loại bỏ các chuỗi
nhiễu trước khi hiển thị. Tesseract là bộ máy OCR mã nguồn mở; tài liệu kỹ
thuật của Tesseract mô tả quy trình nhận dạng văn bản và các thành phần xử lý
ảnh~\cite{Smith2007Tesseract,TesseractDocumentation2026}.

\subsection{Nhận dạng giọng nói (Speech-to-Text)}

Speech-to-Text \en{(STT)} chuyển giọng nói thành văn bản để hệ thống có thể
xử lý câu hỏi của người dùng. Phiên bản hiện tại dùng thư viện
\texttt{SpeechRecognition} gọi Google Web Speech để nhận dạng tiếng Việt.
Luồng này cần kết nối Internet; âm thanh chỉ được gửi sau khi người dùng chủ
động ghi âm và xác nhận xử lý. Source không tích hợp Faster-Whisper hoặc STT
chạy cục bộ~\cite{SpeechRecognitionPackage2026}.

\subsection{Tổng hợp giọng nói (Text-to-Speech)}

Text-to-Speech \en{(TTS)} chuyển phản hồi văn bản thành âm thanh. Hệ thống
ưu tiên Edge TTS với giọng tiếng Việt \texttt{vi-VN-HoaiMyNeural}~\cite{EdgeTTS2026}. Nếu Edge TTS gặp
lỗi, ứng dụng chuyển sang Windows SAPI5 thông qua \texttt{pyttsx3}. File âm thanh
được tạo tạm thời và phát bằng thành phần audio của Streamlit.

\subsection{Khả năng tiếp cận cho người khiếm thị}

Hệ thống hướng tới việc giảm phụ thuộc vào quan sát màn hình bằng cách cho
phép người dùng tải/chụp ảnh, đặt câu hỏi bằng văn bản hoặc giọng nói và
nghe kết quả. Các chức năng hỗ trợ gồm mô tả ảnh, đọc chữ, nhận diện tiền,
phát hiện nguy cơ, hỗ trợ di chuyển và đọc kết quả bằng giọng nói. Các tiêu
chí về focus bàn phím, tương phản và giảm chuyển động có liên quan đến hướng
dẫn tiếp cận nội dung web WCAG~\cite{W3CWCAG22}. Hướng dẫn WAI về văn bản thay
thế cũng phân biệt ảnh thông tin, ảnh trang trí và ảnh chức năng
~\cite{W3CImagesTutorial2026}.

\section{Công cụ và công nghệ}

Bảng dưới đây được tổng hợp trực tiếp từ \texttt{requirements.txt}, \texttt{.env.example}
và cấu trúc mã nguồn của phiên bản VisionAI Pro.

\begin{table}[H]
\centering
\caption{Các công cụ, thư viện và công nghệ được sử dụng trong đồ án.}
\label{tab:congcu}
\begin{tabular*}{\textwidth}{@{\extracolsep{\fill}}>{\bfseries\color{daunavy}}p{3.1cm} >{\color{daumaroon}}p{3.1cm} p{7.0cm}}
\toprule
\rowcolor{gray100}
Thành phần & Phiên bản khai báo tối thiểu & Vai trò trong hệ thống \\
\midrule
Python & Python 3.x & Ngôn ngữ lập trình chính cho giao diện, xử lý ảnh, AI và giọng nói. \\
Streamlit & $\ge$ 1.50 & Framework xây dựng giao diện web chạy cục bộ. \\
Google GenAI SDK & $\ge$ 1.30.0 & Kết nối Gemini để phân tích ảnh và hỏi đáp theo ảnh. \\
Gemini Vision & ID mặc định trong source: \texttt{gemini-3.8-flash} & Đã kiểm tra kết nối bằng \texttt{gemini-3.5-flash-lite}; các model còn lại và độ chính xác chưa được đánh giá. \\
Pillow & $\ge$ 10.0 & Đọc và xử lý ảnh đầu vào. \\
Tesseract + pytesseract & pytesseract $\ge$ 0.3.13 & OCR cục bộ, ưu tiên dữ liệu tiếng Việt và tiếng Anh. \\
SpeechRecognition & $\ge$ 3.10 & STT tiếng Việt thông qua Google Web Speech; cần Internet. \\
Edge TTS & $\ge$ 7.0 & TTS tiếng Việt, ưu tiên giọng \texttt{vi-VN-HoaiMyNeural}. \\
pyttsx3 & $\ge$ 2.98 & TTS dự phòng bằng Windows SAPI5. \\
python-dotenv & $\ge$ 1.0 & Đọc biến môi trường và API key từ file \texttt{.env}. \\
JSON & Có sẵn trong Python & Lưu lịch sử phân tích tại \texttt{data/history.json}, tối đa 100 bản ghi. \\
\bottomrule
\end{tabular*}
\end{table}

Bảng~\ref{tab:congcu} cho biết thư viện nào đảm nhiệm từng công đoạn, nhờ đó
việc bảo trì và kiểm thử có thể tập trung vào từng thành phần. Tách riêng OCR,
STT và TTS cũng giúp chọn dịch vụ phù hợp cho từng loại đầu vào; riêng OCR
chạy cục bộ nên không phải tải ảnh lên Gemini khi người dùng chỉ cần đọc chữ.
Ứng dụng giao diện được xây dựng trên Streamlit~\cite{StreamlitDocumentation2026}.
Các thao tác lấy ảnh từ webcam và ghi âm dựa trên widget
\texttt{st.camera\_input} và \texttt{st.audio\_input} của Streamlit
~\cite{StreamlitCameraInput2026,StreamlitAudioInput2026}.

Không sử dụng cơ sở dữ liệu quan hệ trong phiên bản source hiện tại. Lịch sử
được lưu bằng tệp JSON cục bộ; cấu hình giọng nói, mức chi tiết và trạng thái
hội thoại được quản lý bằng \texttt{st.session\_state} trong phiên Streamlit.

\section{Mô hình AI và cơ chế dự phòng}

Danh sách ID model Gemini được cấu hình theo thứ tự trong biến
\texttt{GEMINI\_MODELS}; giá trị mặc định trong source là
\texttt{gemini-3.8-flash,gemini-3.6-flash,gemini-3.5-flash-lite}. Ứng dụng lần lượt
thử các model khi một yêu cầu gặp lỗi, dừng khi API từ chối khoá/quyền truy cập
(401/403), và chờ ngắn trước khi chuyển model khi gặp giới hạn tốc độ (429).
Danh sách dự phòng không khắc phục được khoá sai, hết hạn mức toàn tài khoản
hoặc lỗi mạng; khả năng truy cập đã được xác nhận với
\texttt{gemini-3.5-flash-lite}, còn các ID khác trong danh sách chưa được kiểm
tra bằng yêu cầu API thật. Cài đặt có nút gửi một yêu cầu nhỏ để kiểm tra
kết nối Gemini.

Yêu cầu Vision sử dụng JSON response mode. Source không khai báo công cụ
function calling; kết quả được chuẩn hoá về một cấu trúc thống nhất trước khi
giao diện hiển thị.

\section{Định hướng lựa chọn công nghệ}

Source tách các tác vụ theo khả năng và phụ thuộc của từng thành phần: Gemini
xử lý mô tả ảnh và hỏi đáp; Tesseract cung cấp OCR cục bộ; SpeechRecognition
chuyển bản ghi thành văn bản qua dịch vụ trực tuyến; Edge TTS và Windows SAPI5
tạo âm thanh phản hồi. Cách phân tách này cho phép dùng OCR khi không cần gọi
Gemini, nhưng không đồng nghĩa toàn bộ ứng dụng chạy ngoại tuyến. Đây là mô tả
kiến trúc source, không phải kết quả so sánh định lượng với hệ thống khác.

\section{Tóm tắt chương}

Chương này đã xác định các công nghệ thực tế của hệ thống gồm Python,
Streamlit, Gemini Vision thông qua Google GenAI SDK, Pillow, Tesseract OCR,
SpeechRecognition/Google Web Speech, Edge TTS, pyttsx3, python-dotenv và JSON.
Các thành phần này tạo thành chuỗi xử lý đa phương thức từ ảnh/giọng nói đến
phân tích AI, văn bản và âm thanh phản hồi.

% ---------------------------------------------------------------------
%  CHƯƠNG 3 — PHÂN TÍCH VÀ THIẾT KẾ
% ---------------------------------------------------------------------
\chapter{Phân tích và thiết kế}
\label{ch:thietke}

\begin{muctieu}
Chương này phân tích các yêu cầu và mô tả kiến trúc thực tế của phiên bản
VisionAI Pro. Trọng tâm là luồng xử lý ảnh, OCR local, Gemini Vision,
STT/TTS, robot Vision, lệnh thoại, hỏi đáp theo ảnh và lưu lịch sử.
\end{muctieu}

\section{Phân tích bài toán và yêu cầu}

Hệ thống được xây dựng nhằm hỗ trợ người khiếm thị tiếp nhận thông tin từ
hình ảnh và tương tác với trợ lý bằng tiếng Việt. Người dùng có thể tải ảnh,
chụp ảnh chủ động bằng camera, nhập câu hỏi hoặc ghi âm rồi xác nhận gửi.
Ảnh được phân tích bằng Gemini Vision hoặc Tesseract OCR tuỳ chế độ; giọng nói
được chuyển thành văn bản bằng Google Web Speech trước khi gửi câu hỏi tới Gemini.

\subsection{Yêu cầu chức năng}

Các chức năng đã được hiện thực trong source gồm:
\begin{itemize}[leftmargin=1.6em, itemsep=2pt]
  \item Hiển thị giao diện Streamlit với 11 trang: Trang chính, Phân tích hình ảnh, OCR, Hỗ trợ di chuyển, Hỏi đáp AI, Trợ lý giọng nói, Camera, Lịch sử, Thống kê, Cài đặt và Hướng dẫn.
  \item Cho phép tải ảnh JPG/JPEG/PNG/WEBP hoặc chụp ảnh bằng camera.
  \item Phân tích ảnh theo các chế độ: mô tả toàn cảnh, nhận biết vật thể và vị trí, OCR, nhận diện tiền Việt, hỗ trợ người khiếm thị và hỗ trợ di chuyển.
  \item Trả về mô tả, bối cảnh, vật thể, nguy cơ, hướng dẫn di chuyển, văn bản và mệnh giá tiền theo cấu trúc dữ liệu chuẩn hoá.
  \item OCR local bằng Tesseract, không sử dụng credit AI cloud.
  \item Cho phép hỏi tiếp về ảnh gần nhất và duy trì tối đa các lượt hội thoại gần nhất trong phiên.
  \item Nhận giọng nói tiếng Việt bằng SpeechRecognition/Google Web Speech; chỉ xử lý bản ghi sau thao tác xác nhận của người dùng.
  \item Chuyển kết quả văn bản thành giọng nói bằng Edge TTS hoặc Windows SAPI5 qua pyttsx3; hỗ trợ lựa chọn giọng đọc.
  \item Cung cấp robot Vision tương tác, hỏi đáp theo ngữ cảnh trang hiện tại và ảnh gần nhất, tìm vật trong ảnh và lệnh thoại điều hướng/đọc lại.
  \item Lưu lịch sử phân tích vào \texttt{data/history.json}; hỗ trợ lọc, đánh dấu yêu thích, xoá có xác nhận, thống kê và xuất báo cáo TXT.
  \item Cung cấp chế độ trợ năng, bật/tắt giọng đọc, mức chi tiết và lựa chọn giọng trong phiên Streamlit.
\end{itemize}

\subsection{Yêu cầu phi chức năng}

\begin{itemize}[leftmargin=1.6em, itemsep=2pt]
  \item \textbf{Khả năng tiếp cận:} ưu tiên phản hồi âm thanh và nội dung ngắn, dễ nghe.
  \item \textbf{Tính ổn định:} kiểm tra ảnh đầu vào, giới hạn dung lượng, báo lỗi dễ hiểu và thử danh sách model Gemini đã cấu hình.
  \item \textbf{An toàn thông tin:} API key được đọc từ \texttt{.env} và không đưa vào source công khai.
  \item \textbf{Khả năng chạy cục bộ:} giao diện chạy bằng Streamlit trên Windows; OCR dùng Tesseract cục bộ, còn Gemini, Google Web Speech và Edge TTS cần mạng.
  \item \textbf{Khả năng mở rộng:} danh sách model Gemini, ngôn ngữ STT và giọng TTS có thể cấu hình qua biến môi trường.
\end{itemize}

\section{Thiết kế kiến trúc hệ thống}

Kiến trúc thực tế được tổ chức thành bốn nhóm: giao diện Streamlit, xử lý
ảnh/âm thanh, các dịch vụ AI và lớp lưu trữ cục bộ. \texttt{app.py} điều phối giao
diện và gọi các module trong thư mục \texttt{modules/}.

\begin{figure}[H]
\centering
\resizebox{\textwidth}{!}{%
\begin{tikzpicture}[
    font=\small\sffamily,
    box/.style={rectangle, rounded corners=5pt, draw=daunavy, fill=daunavy!7, thick, minimum width=4.2cm, minimum height=0.8cm, align=center},
    ai/.style={rectangle, rounded corners=5pt, draw=daumaroon, fill=daumaroon!7, thick, minimum width=4.2cm, minimum height=0.8cm, align=center},
    data/.style={rectangle, rounded corners=5pt, draw=daunavy, fill=white, thick, minimum width=4.2cm, minimum height=0.8cm, align=center},
    arrow/.style={-Stealth, thick, draw=daunavy}
]
\node (user)     [box]  at (1.25,0)    {Người dùng};
\node (ui)       [box]  at (1.25,-1.4) {Streamlit \texttt{app.py}};
\node (image)    [box]  at (-5,-2.9)   {Ảnh tải lên / camera};
\node (question) [box]  at (5,-2.9)    {Câu hỏi văn bản / audio};

\node (vision)   [ai]   at (-7.5,-4.8) {Gemini Vision\\Google GenAI SDK};
\node (ocr)      [ai]   at (-2.5,-4.8) {Tesseract OCR\\local / chế độ OCR};
\node (stt)      [ai]   at (2.5,-4.8)  {SpeechRecognition\\Google Web Speech};
\node (qa)       [ai]   at (7.5,-4.8)  {Gemini hỏi đáp\\ảnh gần nhất nếu có};

\node (analysis) [box]  at (-5,-6.8)   {Kết quả phân tích ảnh};
\node (answer)   [box]  at (7.5,-6.8)  {Trả lời hỏi đáp};
\node (display)  [box]  at (1.25,-8.6) {Hiển thị trong Streamlit};

\node (history)  [data] at (-9,-8.6)   {history.json\\chỉ lưu kết quả\\phân tích};
\node (tts)      [ai]   at (1.25,-10.4) {Edge TTS\\Windows SAPI5 dự phòng\\khi người dùng yêu cầu đọc};
\node (audio)    [box]  at (1.25,-12.2) {Audio tùy chọn};

\draw [arrow] (user) -- (ui);
\draw [arrow] (ui.south) -- ++(0,-0.4) -| (image.north);
\draw [arrow] (ui.south) -- ++(0,-0.4) -| (question.north);
\draw [arrow] (image.south) -- ++(0,-0.4) -| (vision.north);
\draw [arrow] (image.south) -- ++(0,-0.4) -| (ocr.north);
\draw [arrow] (question.south) -- ++(0,-0.4) -| (stt.north);
\draw [arrow] (question.south) -- ++(0,-0.4) -| (qa.north);
\draw [arrow] (stt.east) -- (qa.west);

\draw [arrow] (vision.south) -- ++(0,-0.5) -| (analysis.north);
\draw [arrow] (ocr.south) -- ++(0,-0.5) -| (analysis.north);
\draw [arrow] (qa.south) -- (answer.north);

\draw [arrow] (analysis.west) -| (history.north);
\draw [arrow] (analysis.south) |- (display.west);
\draw [arrow] (answer.south) |- (display.east);
\draw [arrow] (display) -- (tts);
\draw [arrow] (tts) -- (audio);
\end{tikzpicture}
}
\caption{Kiến trúc thực tế của hệ thống Trợ lý thị giác AI.}
\label{fig:kientruc}
\end{figure}

Hình~\ref{fig:kientruc} tách luồng phân tích ảnh khỏi luồng hỏi đáp. Ảnh được
đưa tới Gemini hoặc Tesseract tuỳ chế độ; khi người dùng đặt câu hỏi bằng
giọng nói, STT chuyển bản ghi thành văn bản trước khi gửi cho Gemini. Câu hỏi
nhập bằng bàn phím bỏ qua STT. TTS chỉ được gọi khi người dùng yêu cầu nghe;
kết quả phân tích ảnh được lưu vào lịch sử cục bộ.

Trong source, OCR local là nhánh riêng với Gemini Vision. STT là một nhánh
riêng nhưng gọi Google Web Speech thông qua
\texttt{SpeechRecognition}, nên cần Internet. Kết quả từ Gemini và OCR đều được
chuẩn hoá để giao diện hiển thị; hỏi đáp có thể kèm ảnh gần nhất làm ngữ cảnh.

\section{Thiết kế luồng xử lý hình ảnh}

Luồng xử lý dưới đây bám sát hàm \texttt{analyze\_image()} và \texttt{analyze\_local\_ocr()}
trong source.

\begin{figure}[H]
\centering
\resizebox{0.95\textwidth}{!}{%
\begin{tikzpicture}[
    font=\small\sffamily,
    block/.style={rectangle, rounded corners=5pt, draw=daunavy, fill=white, thick, minimum width=6.2cm, minimum height=0.8cm, align=center},
    ai/.style={rectangle, rounded corners=5pt, draw=daumaroon, fill=daumaroon!7, thick, minimum width=6.2cm, minimum height=0.8cm, align=center},
    decision/.style={diamond, draw=daunavy, fill=daunavy!7, thick, aspect=2.6, align=center, inner sep=1pt},
    arrow/.style={-Stealth, thick, draw=daunavy},
    lab/.style={font=\scriptsize\sffamily, fill=white, inner sep=1pt}
]
\node (start)    [block]    at (0,0)        {Tải ảnh / Chụp camera};
\node (mode)     [decision] at (0,-1.7)     {Chế độ OCR?};
\node (ocr)      [ai]       at (-4.4,-3.9)  {Chọn OCR:\\Tesseract local\\tiền xử lý + nhận dạng};
\node (vision)   [ai]       at (4.4,-3.9)   {Chế độ khác:\\Gemini Vision\\JSON response};
\node (retry)    [block]    at (4.4,-5.9)   {Thử model theo danh sách\\429: chờ ngắn; 401/403: dừng};
\node (fallback) [block]    at (4.4,-8.0)   {OCR local dự phòng\\nếu tất cả model thất bại\\(chỉ đọc chữ)};
\node (normalize)[block]    at (0,-10.2)    {Chuẩn hoá kết quả};
\node (display)  [block]    at (0,-11.6)    {Hiển thị kết quả + cảnh báo};
\node (save)     [block]    at (-4.4,-13.3) {Lưu history.json\\sau phân tích ảnh};
\node (optional) [block]    at (4.4,-13.3)  {TTS / xuất TXT\\khi người dùng yêu cầu};

\draw [arrow] (start) -- (mode);
\draw [arrow] (mode.west) -| (ocr.north) node[lab, pos=0.25] {Có};
\draw [arrow] (mode.east) -| (vision.north) node[lab, pos=0.25] {Không};
\draw [arrow] (vision) -- (retry);
\draw [arrow] (retry) -- (fallback) node[lab, midway] {thất bại};
\draw [arrow] (ocr.south) |- (normalize.west);
\draw [arrow] (retry.west) -| (normalize.north) node[lab, pos=0.25] {thành công};
\draw [arrow] (fallback.south) |- (normalize.east);
\draw [arrow] (normalize) -- (display);
\draw [arrow] (display.south) -- ++(0,-0.4) -| (save.north);
\draw [arrow] (display.south) -- ++(0,-0.4) -| (optional.north);
\end{tikzpicture}
}
\caption{Sơ đồ luồng xử lý ảnh theo source VisionAI Pro.}
\label{fig:luong_xuly_chitiet}
\end{figure}

Các bước chính gồm:
\begin{enumerate}[leftmargin=1.8em, itemsep=3pt]
  \item Người dùng tải ảnh hoặc chụp ảnh từ camera.
  \item Ứng dụng kiểm tra chế độ. Nếu chọn OCR, ảnh được đưa vào Tesseract local.
  \item Với các chế độ Vision, ảnh được mã hoá thành data URL, tách bytes và MIME type rồi gửi tới Gemini.
  \item Gemini được yêu cầu trả về một JSON object gồm mô tả, độ tin cậy, bối cảnh, vật thể, nguy cơ, navigation, OCR và tiền.
  \item Nếu một model Gemini gặp lỗi, ứng dụng lần lượt thử model tiếp theo; với 429, ứng dụng chờ ngắn trước lần thử tiếp theo, còn 401/403 thì dừng.
  \item Nếu toàn bộ yêu cầu Vision thất bại, source thử OCR local và ghi cảnh báo rõ rằng kết quả OCR không thay thế phân tích cảnh/vật thể.
  \item Kết quả được chuẩn hoá, hiển thị trên giao diện và có thể đọc bằng TTS, lưu lịch sử hoặc xuất TXT.
\end{enumerate}

Giải thuật~\ref{alg:xuly-anh} tóm tắt luồng xử lý trên. Nhánh OCR được chọn
trực tiếp khi người dùng chọn chế độ OCR; OCR cũng được thử làm phương án dự
phòng khi mọi yêu cầu Gemini phát sinh lỗi. Nếu OCR dự phòng không khả dụng,
ứng dụng báo lỗi thay vì tạo kết quả phân tích thay thế.

\begin{algorithm}[H]
\caption{Xử lý ảnh theo chế độ OCR hoặc Vision}
\label{alg:xuly-anh}
\hspace*{\algorithmicindent}\textbf{Đầu vào:} ảnh $I$, chế độ xử lý $m$, danh sách model Gemini $M$\\
\hspace*{\algorithmicindent}\textbf{Đầu ra:} kết quả đã chuẩn hoá và cảnh báo (nếu có)
\begin{algorithmic}[1]
\Function{XuLyAnh}{$I, m, M$}
  \If{$m = \text{OCR}$}
    \State $r \gets \Call{OCRLocal}{I}$
  \Else
    \State $r \gets \Call{GeminiVision}{I, M}$
    \If{Tất cả yêu cầu Gemini phát sinh lỗi}
      \State $r \gets \Call{OCRLocal}{I}$
      \State Gắn cảnh báo rằng OCR không thay thế phân tích cảnh và vật thể
    \EndIf
  \EndIf
  \State $r \gets \Call{ChuanHoa}{r}$
  \State Hiển thị $r$ trên giao diện
  \State Lưu kết quả phân tích ảnh vào history.json
  \State TTS và xuất TXT chỉ thực hiện khi người dùng yêu cầu
  \State \Return $r$
\EndFunction
\end{algorithmic}
\end{algorithm}

\section{Thiết kế luồng hỏi đáp và giọng nói}

\begin{figure}[H]
\centering
\resizebox{0.95\textwidth}{!}{%
\begin{tikzpicture}[
    font=\small\sffamily,
    block/.style={rectangle, rounded corners=5pt, draw=daunavy, fill=white, thick, minimum width=5.2cm, minimum height=0.8cm, align=center},
    ai/.style={rectangle, rounded corners=5pt, draw=daumaroon, fill=daumaroon!7, thick, minimum width=5.2cm, minimum height=0.8cm, align=center},
    decision/.style={diamond, draw=daunavy, fill=daunavy!7, thick, aspect=2.4, align=center, inner sep=1pt},
    arrow/.style={-Stealth, thick, draw=daunavy},
    lab/.style={font=\scriptsize\sffamily, fill=white, inner sep=1pt}
]
\node (text)        [block]    at (-3.6,0)    {Câu hỏi nhập bằng bàn phím};
\node (audio_input) [block]    at (3.6,0)     {Ghi âm + xác nhận\\của người dùng};
\node (stt)         [ai]       at (3.6,-1.9)  {SpeechRecognition\\Google Web Speech};
\node (parser)      [decision] at (0,-4.0)    {Lệnh cục bộ?};
\node (local)       [block]    at (-4.4,-6.4) {Lệnh cục bộ:\\điều hướng / trợ giúp\\đọc lại};
\node (qa)          [ai]       at (4.4,-6.4)  {Hỏi Gemini:\\tìm vật trong ảnh\\ảnh gần nhất nếu có};
\node (answer)      [block]    at (4.4,-8.2)  {Câu trả lời bằng văn bản};
\node (tts)         [ai]       at (4.4,-10.0) {Edge TTS\\Windows SAPI5 dự phòng\\nếu người dùng yêu cầu đọc};
\node (player)      [block]    at (4.4,-11.8) {Audio tùy chọn};

\draw [arrow] (text.south) |- (parser.west);
\draw [arrow] (audio_input) -- (stt);
\draw [arrow] (stt.south) |- (parser.east);
\draw [arrow] (parser.south west) -| (local.north) node[lab, pos=0.25] {Có};
\draw [arrow] (parser.south east) -| (qa.north) node[lab, pos=0.25] {Không};
\draw [arrow] (qa) -- (answer);
\draw [arrow] (answer) -- (tts);
\draw [arrow] (tts) -- (player);
\end{tikzpicture}
}
\caption{Luồng xử lý giọng nói, hỏi đáp theo ảnh và phản hồi âm thanh.}
\label{fig:luong_giongnoi}
\end{figure}

Hình~\ref{fig:luong_giongnoi} mô tả trình tự sau khi người dùng chủ động
ghi âm và xác nhận: nhận dạng trực tuyến, xử lý lệnh thoại hoặc gửi câu hỏi
cùng ngữ cảnh tới Gemini, rồi phát câu trả lời bằng TTS khi được yêu cầu.
Các mũi tên thể hiện thứ tự xử lý, không phải luồng âm thanh thời gian thực.

Người dùng chủ động ghi âm và xác nhận để \texttt{SpeechRecognition} gửi âm
thanh tới Google Web Speech chuyển thành văn bản. Câu hỏi sau đó được gửi cùng
ngữ cảnh trang và ảnh gần nhất (nếu có) tới Gemini. Kết quả có thể được đọc
bằng Edge TTS hoặc Windows SAPI5 khi Edge TTS gặp lỗi. Luồng này không phải
hội thoại âm thanh thời gian thực.

\section{Thiết kế cơ chế cá nhân hoá}

Phiên bản source hiện tại thực hiện cá nhân hoá ở mức phiên làm việc thay vì
xây dựng hồ sơ người dùng hoặc cơ sở dữ liệu. \texttt{st.session\_state} lưu trạng thái
trợ năng, bật/tắt giọng nói, giọng đọc, mức chi tiết, chế độ hiện tại, ảnh gần nhất,
ngữ cảnh trang và lịch sử chat. Ngoài ra, \texttt{history.json} lưu tối đa 100 bản ghi
phân tích để người dùng xem lại, lọc và đánh dấu yêu thích.

\begin{figure}[H]
\centering
\begin{tikzpicture}[
    node distance=0.85cm and 0.9cm,
    font=\small\sffamily,
    box/.style={rectangle, rounded corners=5pt, draw=daunavy, fill=white, thick, minimum width=3.3cm, minimum height=0.72cm, align=center},
    center/.style={rectangle, rounded corners=5pt, draw=daumaroon, fill=daumaroon!7, thick, minimum width=3.6cm, minimum height=0.9cm, align=center},
    arrow/.style={-Stealth, thick, draw=daunavy}
]
\node (settings) [box] {Voice / Detail / Mode};
\node (image) [box, below=of settings] {Ảnh gần nhất};
\node (chat) [box, below=of image] {Ngữ cảnh chat};
\node (session) [center, right=1.3cm of image] {st.session\_state};
\node (analysis) [box, right=1.3cm of session] {Kết quả phân tích ảnh};
\node (history) [box, below=of analysis] {history.json\\tối đa 100 bản ghi};
\node (result) [box, below=1.0cm of session] {Kết quả và trải nghiệm phiên};
\draw [arrow] (settings) -- (session);
\draw [arrow] (image) -- (session);
\draw [arrow] (chat) -- (session);
\draw [arrow] (session) -- (result);
\draw [arrow] (session) -- (analysis);
\draw [arrow] (analysis) -- (history);
\end{tikzpicture}
\caption{Cơ chế cá nhân hoá và lưu trạng thái thực tế của phiên bản hiện tại.}
\label{fig:canhanhoa}
\end{figure}

Hình~\ref{fig:canhanhoa} thể hiện các tuỳ chọn và ngữ cảnh tạm thời được giữ
trong phiên Streamlit, còn các kết quả được lưu riêng trong tệp lịch sử.
Điều này giải thích vì sao phiên bản hiện tại chưa có hồ sơ cá nhân lâu dài.

Cách tiếp cận này phù hợp với phạm vi nguyên mẫu. Source chưa có đăng nhập,
hồ sơ người dùng hoặc cơ sở dữ liệu quan hệ; vì vậy báo cáo không mô tả các
thành phần này như một chức năng đã triển khai.

\section{Thiết kế dữ liệu}

Dữ liệu phân tích được biểu diễn bằng cấu trúc JSON gồm các trường \texttt{summary},
\texttt{confidence}, \texttt{scene}, \texttt{objects}, \texttt{accessibility\_tips}, \texttt{text\_detected},
\texttt{money\_detected}, \texttt{hazards}, \texttt{navigation} và \texttt{warning}. Lịch sử được lưu trong
\texttt{data/history.json} và mỗi lần lưu chỉ giữ tối đa 100 bản ghi gần nhất.

\section{Thiết kế giao diện}

Giao diện được xây dựng bằng Streamlit với thanh điều hướng bên trái và một
robot Vision đồ hoạ có thể mở chat hoặc điều hướng tới các công cụ. Các trang
gồm trang chính, phân tích ảnh, OCR, hỗ trợ di chuyển, hỏi đáp AI, trợ lý giọng
nói, camera, lịch sử, thống kê, cài đặt và hướng dẫn. Chế độ trợ năng tăng cỡ
chữ/nút và làm rõ focus bàn phím; người dùng có thể đổi giọng đọc, bật/tắt
giọng nói, kiểm tra kết nối Gemini và xoá phiên hoặc lịch sử.

Hình~\ref{fig:bocuc-web} mô tả bố cục trang chính. Thanh bên chứa điều hướng
11 trang cùng các công tắc trợ năng; khu vực chính gồm robot Vision, hàng
thẻ công cụ, kết quả gần nhất và cột hoạt động gần đây.

\begin{figure}[H]
\centering
\resizebox{\textwidth}{!}{%
\begin{tikzpicture}[
    font=\scriptsize\sffamily,
    frame/.style={rectangle, rounded corners=4pt, draw=daunavy, thick, fill=white},
    side/.style={rectangle, rounded corners=4pt, draw=daunavy, thick, fill=daunavy!7},
    hero/.style={rectangle, rounded corners=4pt, draw=daumaroon, thick, fill=daumaroon!7},
    card/.style={rectangle, rounded corners=3pt, draw=daunavy, thick, fill=daunavy!7, align=center, minimum width=1.38cm, minimum height=1.3cm, inner sep=2pt},
    btn/.style={rectangle, rounded corners=2pt, draw=daunavy, fill=white, align=center, inner sep=2pt, minimum height=0.42cm},
    lbl/.style={font=\scriptsize\sffamily\bfseries, text=daunavy, anchor=west}
]
% khung ngoài
\draw[frame] (0,0) rectangle (15.4,8.9);
% ---------- thanh bên ----------
\draw[side] (0.15,0.15) rectangle (3.75,8.75);
\node[anchor=north west, font=\scriptsize\sffamily\bfseries, text=daumaroon] at (0.3,8.6) {VisionAI Pro};
\node[anchor=north west, font=\tiny\sffamily] at (0.3,8.25) {Trạng thái AI: Đã sẵn sàng};
\node[anchor=north west, align=left, font=\tiny\sffamily, text width=3.2cm] at (0.3,7.85)
  {Điều hướng (11 trang):\\[1pt]
   Trang chính\\ Phân tích hình ảnh\\ OCR\\ Hỗ trợ di chuyển\\ Hỏi đáp AI\\
   Trợ lý giọng nói\\ Camera\\ Lịch sử\\ Thống kê\\ Cài đặt\\ Hướng dẫn};
\node[btn, minimum width=3.2cm] at (1.95,3.05) {Trò chuyện với Vision};
\node[btn, minimum width=3.2cm] at (1.95,2.4) {Công tắc giọng nói};
\node[btn, minimum width=3.2cm] at (1.95,1.75) {Công tắc trợ năng};
\node[btn, minimum width=3.2cm] at (1.95,0.9) {Xóa phiên hiện tại};
% ---------- robot ----------
\draw[hero] (3.95,6.55) rectangle (15.25,8.75);
\node[btn, minimum width=4.2cm, align=left] at (6.15,8.2) {Lời thoại của robot Vision};
\node[anchor=west, font=\small\sffamily\bfseries] at (4.15,7.45) {VisionAI Pro};
\node[btn, anchor=west] (hb1) at (4.15,6.9) {Phân tích hình ảnh};
\node[btn, anchor=west] (hb2) at ([xshift=0.2cm]hb1.east) {Đọc văn bản};
\node[btn, anchor=west] at ([xshift=0.2cm]hb2.east) {Hỏi đáp theo ảnh};
\node[rectangle, rounded corners=6pt, draw=daumaroon, fill=white, minimum width=2.4cm, minimum height=1.7cm, align=center] at (13.3,7.65) {Robot Vision\\(16 trạng thái\\biểu cảm)};
% ---------- cột chính ----------
\node[lbl] at (3.95,6.2) {Công cụ};
\foreach \i/\t in {0/{Phân tích\\ảnh},1/{Đọc\\văn bản},2/{Hỏi đáp\\ảnh},3/{Giọng\\nói},4/{Camera}} {
  \node[card] at (4.65+\i*1.52,5.35) {\t};
  \node[btn, minimum width=1.38cm, font=\tiny\sffamily] at (4.65+\i*1.52,4.45) {Mở};
}
\node[lbl] at (3.95,3.85) {Kết quả gần nhất};
\draw[frame] (3.95,0.3) rectangle (11.55,3.6);
\node[align=left, anchor=north west, font=\tiny\sffamily] at (4.1,3.45)
  {Mô tả, bối cảnh, cảnh báo an toàn, vật thể,\\
   tiền Việt, văn bản OCR, gợi ý cho người khiếm thị,\\
   nút nghe kết quả và nút xuất báo cáo TXT.};
% ---------- cột phải ----------
\node[btn, minimum width=3.4cm] at (13.4,6.2) {Cấu hình Gemini};
\draw[frame] (11.75,0.3) rectangle (15.25,5.85);
\node[anchor=north west, font=\scriptsize\sffamily\bfseries, text=daunavy] at (11.9,5.7) {Hoạt động gần đây};
\node[anchor=north west, align=left, font=\tiny\sffamily] at (11.9,5.25)
  {Tối đa 4 bản ghi\\từ \texttt{history.json}\\(chế độ và thời điểm\\phân tích).};
\end{tikzpicture}}
\caption[Sơ đồ bố cục trang chính của giao diện web]{Sơ đồ bố cục trang chính của giao diện web, vẽ lại từ cấu trúc trong \texttt{app.py}; đây là sơ đồ thiết kế, không phải ảnh chụp màn hình.}
\label{fig:bocuc-web}
\end{figure}

\section{Thiết kế thành phần AI}

Thành phần AI được tách thành \texttt{modules/ai\_vision.py}, \texttt{modules/local\_ocr.py}
và \texttt{modules/speech.py}. \texttt{ai\_vision.py} phụ trách Gemini Vision và hỏi đáp theo
ảnh; \texttt{local\_ocr.py} phụ trách Tesseract; \texttt{speech.py} phụ trách STT/TTS. Cách
tách module giúp thay đổi từng thành phần mà không phải sửa toàn bộ giao diện.
%% ---------------------------------------------------------------------
%  CHƯƠNG 4 — CÀI ĐẶT VÀ THỬ NGHIỆM
% ---------------------------------------------------------------------
\chapter{Cài đặt và thử nghiệm}
\label{ch:caidat}

\begin{muctieu}
Chương này trình bày môi trường phát triển và triển khai, cấu trúc mã nguồn,
các thành phần đã cài đặt và quy trình chạy thực tế của phiên bản
VisionAI Pro.
\end{muctieu}

\section{Môi trường phát triển và triển khai}

Theo tài liệu \texttt{README.md} của source, hệ thống được hướng dẫn chạy trên
Windows bằng PowerShell hoặc terminal tại thư mục chứa \texttt{app.py}. Ứng dụng
được triển khai theo hình thức chạy cục bộ bằng Streamlit, không có cấu hình
server production hoặc dịch vụ cloud riêng trong source hiện tại.

\begin{table}[H]
\centering
\caption{Môi trường phát triển và triển khai được xác định từ source.}
\label{tab:moitruong}
\begin{tabular*}{\textwidth}{@{\extracolsep{\fill}}>{\bfseries\color{daunavy}}p{3.0cm} >{\color{daumaroon}}p{4.2cm} p{6.0cm}}
\toprule
\rowcolor{gray100}
Thành phần & Thông tin & Ghi chú \\
\midrule
Hệ điều hành & Windows & README hướng dẫn chạy bằng PowerShell; script OCR sử dụng winget và đường dẫn Tesseract của Windows. \\
Python & Python 3.11 hoặc 3.12 & README hướng dẫn hai phiên bản này; requirements không khóa phiên bản Python. \\
Giao diện & Streamlit $\ge$ 1.50 & Chạy bằng lệnh \texttt{streamlit run app.py}. \\
AI Vision & Gemini theo \texttt{GEMINI\_MODELS} & Đã kiểm tra kết nối và chạy ảnh mẫu bằng \texttt{gemini-3.5-flash-lite}; chưa đánh giá độ chính xác tổng quát. \\
OCR & Tesseract + pytesseract & Chạy local; setup script tải \texttt{vie.traineddata} và \texttt{eng.traineddata}. \\
STT & SpeechRecognition / Google Web Speech & Nhận dạng tiếng Việt qua mạng; không có STT Whisper cục bộ. \\
TTS & Edge TTS / Windows SAPI5 & Edge TTS ưu tiên; \texttt{pyttsx3} dùng SAPI5 khi cần. \\
Cấu hình & \texttt{.env} & Lưu API key, danh sách model, ngôn ngữ STT, giọng TTS và đường dẫn Tesseract. \\
Lưu trữ & \texttt{data/history.json} & Lưu lịch sử phân tích cục bộ, tối đa 100 bản ghi. \\
\bottomrule
\end{tabular*}
\end{table}

Bảng~\ref{tab:moitruong} tóm tắt các yêu cầu môi trường lấy từ hướng dẫn chạy
source; các thông số phần cứng không có trong tài liệu nên không được suy đoán.
Không có đủ thông tin từ source để xác định chính xác CPU, RAM hoặc phiên bản
Python đã cài trên máy phát triển. Vì vậy các thông số phần cứng này không
được đưa thành số liệu chính thức trong báo cáo.

\section{Cài đặt và chạy hệ thống}

Quy trình chạy được mô tả trong \texttt{README.md} và có thể thực hiện theo các bước:

\begin{lstlisting}[language=bash, caption={Cài đặt thư viện và khởi chạy ứng dụng trên Windows.}]
pip install -r requirements.txt
streamlit run app.py
\end{lstlisting}

Trước khi chạy chức năng Gemini, người dùng cần tạo file \texttt{.env} và khai báo
\texttt{GEMINI\_API\_KEY}. Không đưa giá trị khoá vào báo cáo hoặc ảnh chụp.
Các biến trong mẫu cấu hình gồm \texttt{GEMINI\_MODELS}, \texttt{EDGE\_TTS\_VOICE},
\texttt{EDGE\_TTS\_RATE}, \texttt{EDGE\_TTS\_PITCH}, \texttt{STT\_LANGUAGE} và
\texttt{TESSERACT\_CMD}.

Đối với OCR tiếng Việt trên Windows, source cung cấp \texttt{setup\_ocr.ps1}. Script
có thể yêu cầu quyền Administrator, sử dụng \texttt{winget} để cài Tesseract và tải
các dữ liệu ngôn ngữ \texttt{vie.traineddata} và \texttt{eng.traineddata}.

\section{Cấu trúc mã nguồn}

Cấu trúc chính của source gồm:

\begin{table}[H]
\centering
\caption{Các tệp và thư mục chính của hệ thống.}
\label{tab:cautrucsource}
\begin{tabular*}{\textwidth}{@{\extracolsep{\fill}}>{\bfseries\color{daunavy}}p{4.1cm} p{9.0cm}}
\toprule
\rowcolor{gray100}
Tệp / thư mục & Chức năng \\
\midrule
\texttt{app.py} & Điểm vào chính, giao diện Streamlit, điều hướng, hiển thị kết quả và điều phối các module. \\
\texttt{modules/ai\_vision.py} & Gemini Vision, hỏi đáp theo ảnh, chuẩn hoá JSON và thử model theo danh sách. \\
\texttt{modules/local\_ocr.py} & Tiền xử lý ảnh và OCR local bằng Tesseract. \\
\texttt{modules/speech.py} & STT Google Web Speech và TTS Edge/Windows. \\
\texttt{modules/robot\_ui.py}, \texttt{ui/vision\_robot/} & Robot Vision tương tác và các lối tắt điều hướng. \\
\texttt{modules/accessibility.py} & Tạo nội dung đọc thành tiếng và xử lý văn bản OCR. \\
\texttt{modules/voice\_commands.py} & Phân tích lệnh thoại điều hướng, tìm vật, trợ giúp và đọc lại. \\
\texttt{modules/utils.py} & Chuẩn hoá ảnh, lưu/đọc lịch sử, yêu thích và xuất báo cáo TXT. \\
\texttt{data/history.json} & Tệp lịch sử cục bộ được tạo khi có kết quả; không đóng gói dữ liệu lịch sử. \\
\texttt{tests/test\_core.py} & Kiểm thử cục bộ cho xử lý ảnh, lịch sử, trợ năng, lệnh thoại và yêu cầu Gemini. \\
\texttt{ui/} & Thành phần giao diện web nhúng cho robot. \\
\texttt{setup\_ocr.ps1} & Cài Tesseract và dữ liệu OCR tiếng Việt/Anh trên Windows. \\
\texttt{.env.example} & Mẫu cấu hình biến môi trường; không chứa API key thật. \\
\texttt{requirements.txt} & Danh sách thư viện Python cần cài đặt. \\
\bottomrule
\end{tabular*}
\end{table}

Bảng~\ref{tab:cautrucsource} mô tả trách nhiệm của từng module và vị trí các
kiểm thử, giúp người đọc lần theo luồng xử lý trong mã nguồn.

Gói source đi kèm báo cáo chứa mã ứng dụng, các module cần thiết, kiểm thử,
README và cấu hình mẫu có trường API key để trống. Gói không chứa \texttt{.env},
môi trường ảo, thư mục Git hoặc lịch sử runtime; khoá API cần được tạo riêng
trên máy chạy ứng dụng và không đưa vào kho công khai.
\section{Các thành phần chính}

\begin{itemize}[leftmargin=1.6em, itemsep=2pt]
  \item \textbf{Giao diện:} Streamlit cung cấp 11 trang, robot Vision, chế độ trợ năng, upload/camera/audio input, kết quả và nút xuất báo cáo.
  \item \textbf{AI Vision:} Gemini phân tích ảnh, tạo JSON có cấu trúc và trả lời câu hỏi theo ảnh gần nhất.
  \item \textbf{OCR local:} Tesseract đọc chữ mà không gọi AI cloud.
  \item \textbf{STT:} SpeechRecognition gửi bản ghi đã xác nhận tới Google Web Speech; cần Internet.
  \item \textbf{TTS:} Edge TTS tạo MP3; Windows SAPI5/pyttsx3 là fallback.
  \item \textbf{Lịch sử và xuất báo cáo:} \texttt{history.json} lưu tối đa 100 bản ghi; \texttt{export\_report()} tạo báo cáo TXT.
\end{itemize}

\begin{figure}[H]
\centering
\includegraphics[width=\textwidth]{visionai-ui.png}
\caption{Ảnh giao diện trang chính VisionAI Pro do người dùng cung cấp, gồm robot Vision và các thẻ công cụ; không phải kết quả phân tích ảnh.}
\label{fig:visionai-ui}
\end{figure}

Hình~\ref{fig:visionai-ui} ghi lại giao diện ứng dụng đang chạy; đây là
ảnh chụp trang chính để minh hoạ cách người dùng mở các công cụ, không phải
ảnh chụp kết quả phân tích Gemini.

\section{Cài đặt các chức năng chính}

\subsection{Cài đặt phân tích hình ảnh}

Ảnh được lấy từ \texttt{st.file\_uploader} hoặc \texttt{st.camera\_input}; camera
chỉ chụp sau thao tác chủ động của người dùng, không truyền video trực tiếp.
Ứng dụng giới hạn ảnh 15 MB, kiểm tra ảnh hợp lệ, áp dụng hướng EXIF, đổi sang
RGB, thu nhỏ cạnh tối đa 1600 pixel và mã hoá JPEG trước khi gửi tới module xử lý.
Các giới hạn này đặt trần cho dung lượng đầu vào, giảm số pixel cần mã hoá
và lượng dữ liệu truyền tới dịch vụ so với ảnh gốc độ phân giải cao. Đây là
lý do thiết kế nhằm giảm tải xử lý/truyền dữ liệu; báo cáo không khẳng định
mức tăng tốc định lượng khi chưa đo thời gian thực tế.

\begin{lstlisting}[language=Python, caption={Điều phối chế độ OCR và Gemini Vision trong app.py.}]
if mode == "Đọc chữ trong ảnh (OCR)":
    result = analyze_local_ocr(raw)
else:
    result = analyze_image(data_url, mode, detail, scene_style)
\end{lstlisting}

\begin{table}[H]
\centering
\caption{Các bước tối ưu đầu vào và giới hạn dữ liệu trong pipeline.}
\label{tab:toiuu-anh}
\begin{tabular*}{\textwidth}{@{\extracolsep{\fill}}>{\bfseries\color{daunavy}}p{3.0cm} p{4.1cm} p{6.0cm}}
\toprule
\rowcolor{gray100}
Bước & Giới hạn trong source & Mục đích / tác động dự kiến \\
\midrule
Kiểm tra đầu vào & Tối đa 15 MB; từ chối tệp rỗng hoặc ảnh lỗi & Giới hạn dữ liệu cần đọc và tránh đưa tệp không hợp lệ vào pipeline. \\
Chuẩn hoá ảnh Vision & Hướng EXIF, RGB, cạnh tối đa 1600 px, JPEG chất lượng 88 & Giảm số pixel và payload gửi tới Gemini so với ảnh gốc độ phân giải cao. \\
Tiền xử lý OCR & Cạnh tối đa 2600 px; ảnh grayscale và nhị phân & Giữ kích thước xử lý hữu hạn, tạo biến thể để Tesseract thử nhận diện chữ. \\
Lịch sử & Tối đa 100 bản ghi; thay thế tệp nguyên tử & Giữ dữ liệu lưu trữ hữu hạn và giảm nguy cơ tệp JSON bị ghi dở. \\
\bottomrule
\end{tabular*}
\end{table}

Bảng~\ref{tab:toiuu-anh} giải thích tác dụng của các giới hạn trong pipeline.
Thu nhỏ và nén ảnh có thể giảm thời gian mã hoá/truyền khi ảnh nguồn lớn,
nhưng chất lượng mạng, dịch vụ Gemini và nội dung ảnh cũng ảnh hưởng thời gian
đáp ứng. Muốn kết luận ứng dụng chạy nhanh hơn bao nhiêu cần đo trên cùng bộ
ảnh, thiết bị và kết nối; hiện chưa có phép đo như vậy.

\subsection{Cài đặt Gemini Vision}

\texttt{ai\_vision.py} tạo client bằng API key trong \texttt{.env}, gửi ảnh cùng prompt tới
Gemini và yêu cầu kết quả JSON. Prompt quy định các trường bắt buộc như
\texttt{summary}, \texttt{objects}, \texttt{hazards}, \texttt{navigation}, \texttt{text\_detected} và
\texttt{money\_detected}. Kết quả sau đó được chuẩn hoá để tránh lỗi dữ liệu trước khi
hiển thị. Nội dung do AI trả về được escape trước khi chèn vào HTML của giao
diện để tránh làm hỏng cấu trúc hiển thị bởi các ký tự đặc biệt.

Ứng dụng lần lượt thử các model trong \texttt{GEMINI\_MODELS}; khi gặp 429,
ứng dụng chờ ngắn trước khi thử model tiếp theo. Lỗi 401/403 khiến quá trình
dừng và báo cấu hình/quyền truy cập, thay vì lặp lại vô hạn. Trang Cài đặt có
nút gửi một yêu cầu nhỏ để kiểm tra kết nối Gemini. Nếu mọi yêu cầu Vision
thất bại, hàm xử lý thử Tesseract local và gắn cảnh báo rằng kết quả thay thế
chỉ là văn bản OCR, không phải mô tả cảnh hay vật thể.

Trong lần chạy thử thực tế, kiểm tra kết nối Gemini thành công với model
\texttt{gemini-3.5-flash-lite}. Chế độ nhận diện tiền mô tả tờ tiền 200.000
đồng với độ tin cậy 0,98; chế độ hỗ trợ di chuyển phát cảnh báo nguy cơ cao
trên ảnh giao thông tĩnh với độ tin cậy 0,90. Chế độ mô tả toàn cảnh và nhận
biết vật thể/vị trí cũng được chạy trên ảnh phong cảnh; chế độ hỗ trợ người
khiếm thị được chạy trên ảnh tiền và trả về ba dòng chữ nhận diện. Đây là các
lần thử chức năng trên ảnh đơn lẻ, không phải phép đo độ chính xác của mô hình.
Hỏi đáp ảnh và tìm vật cũng đã gọi Gemini, nhưng giao diện hiển thị câu trả
lời dạng JSON thô thay vì câu trả lời tự nhiên.

\subsection{Cài đặt OCR local}

\texttt{local\_ocr.py} sử dụng Pillow để chuẩn bị ảnh và pytesseract để gọi Tesseract.
Ảnh được chỉnh hướng EXIF, thu nhỏ giới hạn 2600 pixel, chuyển grayscale,
tăng tương phản/làm nét và tạo hai biến thể xử lý. Tesseract thử PSM 6 và 11,
lọc từ có độ tin cậy dưới 12\%, gom theo dòng và chọn một kết quả để tránh
ghép lặp. Hệ thống dùng ngôn ngữ \texttt{vie}, \texttt{eng} hoặc cả hai tuỳ dữ
liệu đã cài.

Máy kiểm tra chỉ có dữ liệu Tesseract cài đặt hệ thống \texttt{eng} và
\texttt{osd}, thiếu \texttt{vie.traineddata}. Để thử lại mà không thay đổi cài
đặt hệ thống, dữ liệu \texttt{vie} được đặt trong thư mục tạm và truyền qua
\texttt{TESSDATA\_PREFIX}. Trên ảnh tiền, Tesseract hiển thị độ tin cậy tổng
hợp 84\%, nhưng vẫn nhận sai nhiều nội dung: \texttt{CONG:HOA XÃ HỘI CHỦ NGHĨA
tM} thay vì ``CỘNG HÒA XÃ HỘI CHỦ NGHĨA'' và \texttt{Bee: NGHÌN DONG} thay
vì ``HAI TRĂM NGHÌN ĐỒNG''. Ảnh có vùng phản sáng; độ tin cậy do OCR trả về
không đồng nghĩa với độ chính xác ngữ nghĩa. Người dùng cần cài dữ liệu
\texttt{vie} theo script đi kèm và kiểm tra lại nội dung quan trọng.

\subsection{Cài đặt STT}

Người dùng bấm ghi âm, sau đó bấm xác nhận nhận diện trước khi bản ghi được
chuyển thành văn bản. \texttt{speech.py} lưu tạm âm thanh, giới hạn dung lượng
15 MB và gọi \texttt{SpeechRecognition.recognize\_google()} với ngôn ngữ
\texttt{vi-VN}; yêu cầu mạng có timeout 15 giây. File tạm được xoá sau xử lý.
Đây là STT trực tuyến, không phải Whisper chạy trên máy.

\subsection{Cài đặt TTS}

\texttt{text\_to\_speech()} ưu tiên Edge TTS với giọng được chọn trong cài đặt.
Khi Edge TTS không tạo được file, chương trình sử dụng Windows SAPI5 qua pyttsx3. File
âm thanh được tạo tạm thời và đưa vào \texttt{st.audio} để phát cho người dùng.

\subsection{Lịch sử, thống kê và xuất báo cáo}

\texttt{save\_history()} ghi nguyên tử bản ghi phân tích vào
\texttt{data/history.json} và giới hạn lịch sử ở 100 bản ghi gần nhất. Trang Lịch
sử hỗ trợ lọc và đánh dấu yêu thích; xoá toàn bộ yêu cầu xác nhận. Trang Thống
kê tổng hợp dữ liệu đã lưu. Hàm \texttt{export\_report()} tạo file TXT gồm mô tả,
vật thể, cảnh báo, OCR, tiền, gợi ý và điều hướng.

\section{Cơ chế cá nhân hoá}

Source hiện thực cá nhân hoá ở mức phiên làm việc thông qua \texttt{st.session\_state}:
người dùng có thể bật/tắt giọng đọc, chọn mức chi tiết Thấp/Vừa/Cao, chọn giọng
đọc, bật chế độ trợ năng và giữ ảnh/ngữ cảnh hội thoại gần nhất. Đây chưa phải
hệ thống hồ sơ cá nhân có đăng nhập hoặc cơ sở dữ liệu.

\section{Giao diện và các chức năng}

Giao diện có 11 trang: Trang chính, Phân tích hình ảnh, OCR, Hỗ trợ di chuyển,
Hỏi đáp AI, Trợ lý giọng nói, Camera, Lịch sử, Thống kê, Cài đặt và Hướng dẫn.
Robot Vision có thể mở chat hoặc lối tắt chức năng; lệnh thoại hỗ trợ điều
hướng trang, tìm vật trong ảnh gần nhất, trợ giúp và đọc lại. Chế độ trợ năng
tăng cỡ chữ/nút và làm rõ focus bàn phím.

Hình~\ref{fig:banner-visionai} là tài nguyên đồ hoạ thương hiệu được cung cấp
cùng source, không phải ảnh chụp màn hình hay bằng chứng về kết quả phân tích.
Ảnh giao diện trang chính do người dùng cung cấp được trình bày ở
Hình~\ref{fig:visionai-ui}.

\begin{figure}[H]
\centering
\includegraphics[width=\textwidth]{visionai_hero_banner.png}
\caption[Banner nhận diện VisionAI Pro]{Tài nguyên banner đồ hoạ VisionAI Pro (không phải ảnh chụp màn hình).}
\label{fig:banner-visionai}
\end{figure}

Bảng~\ref{tab:congcu-giaodien} tổng hợp 11 công cụ (trang) của giao diện web,
đối chiếu với danh sách điều hướng trong \texttt{app.py}; cột cuối nêu thành
phần mã nguồn phụ trách chính.

\begin{table}[H]
\centering
\small\linespread{1.05}\selectfont
\caption{Các công cụ hỗ trợ trên giao diện web của VisionAI Pro.}
\label{tab:congcu-giaodien}
\begin{tabular*}{\textwidth}{@{\extracolsep{\fill}}>{\raggedright\arraybackslash\bfseries\color{daunavy}}p{2.9cm} >{\raggedright\arraybackslash}p{8.3cm} >{\raggedright\arraybackslash}p{3.3cm}}
\toprule
\rowcolor{gray100}
Công cụ (trang) & Hỗ trợ người dùng & Thành phần phụ trách \\
\midrule
Trang chính & Hiển thị robot Vision, năm thẻ công cụ (phân tích ảnh, đọc văn bản, hỏi đáp ảnh, giọng nói, camera), kết quả gần nhất và hoạt động gần đây. & \texttt{app.py}, \texttt{robot\_ui.py} \\
Phân tích hình ảnh & Tải ảnh hoặc chụp ảnh, chọn một trong sáu chế độ (mô tả toàn cảnh, nhận biết vật thể và vị trí, OCR, nhận diện tiền Việt, hỗ trợ người khiếm thị, hỗ trợ di chuyển) rồi phân tích. & \texttt{ai\_vision.py}, \texttt{local\_ocr.py} \\
OCR & Đọc chữ trong ảnh bằng Tesseract cục bộ; kết quả theo từng dòng và có thể nghe lại. & \texttt{local\_ocr.py}, \texttt{accessibility.py} \\
Hỗ trợ di chuyển & Phân tích ảnh phía trước để ưu tiên lối đi, vật cản và cầu thang; luôn kèm cảnh báo rằng ảnh tĩnh không thay thế gậy hay chó dẫn đường. & \texttt{ai\_vision.py} \\
Hỏi đáp AI & Đặt nhiều câu hỏi liên tiếp về cùng một ảnh, giữ ngữ cảnh hội thoại gần nhất. & \texttt{ai\_vision.py} \\
Trợ lý giọng nói & Ghi âm câu hỏi tiếng Việt, xác nhận nhận diện rồi gửi cho Gemini; hỗ trợ lệnh thoại. & \texttt{speech.py}, \texttt{voice\_commands.py} \\
Camera & Chụp ảnh theo thao tác chủ động của người dùng rồi phân tích; không truyền video trực tiếp. & \texttt{st.camera\_input} \\
Lịch sử & Xem lại các lần phân tích, đánh dấu và lọc mục yêu thích, xoá có xác nhận. & \texttt{utils.py}, \texttt{history.json} \\
Thống kê & Tổng hợp số lần phân tích, số lần có cảnh báo, số lần OCR, nhận diện tiền và phân bố theo chế độ. & \texttt{app.py} \\
Cài đặt & Bật/tắt nút đọc kết quả, chọn mức chi tiết Thấp/Vừa/Cao, chọn giọng đọc, kiểm tra kết nối Gemini. & \texttt{ai\_vision.py}, \texttt{speech.py} \\
Hướng dẫn & Trình bày cách dùng các công cụ và lưu ý an toàn khi sử dụng. & \texttt{app.py} \\
\bottomrule
\end{tabular*}
\end{table}

Robot Vision là thành phần web riêng (\texttt{ui/vision\_robot/index.html}) nhận
trạng thái và lời thoại từ \texttt{app.py}. Hình~\ref{fig:robot-trangthai}
gồm bốn ảnh người dùng cung cấp để minh hoạ trạng thái \texttt{focused},
\texttt{excited}, \texttt{concerned} và \texttt{empathetic}; các trạng thái này có
trong mã nguồn. Ảnh robot minh hoạ trạng thái giao diện, không phải ảnh chụp
một phiên phân tích Gemini cụ thể.

\begin{figure}[H]
\centering
\begin{minipage}[t]{0.485\textwidth}
\centering
\includegraphics[width=\linewidth]{ui-robot-tap-trung.png}\\[3pt]
{\small\textbf{(a)} Tập trung (\texttt{focused}): đang phân tích ảnh}
\end{minipage}\hfill
\begin{minipage}[t]{0.485\textwidth}
\centering
\includegraphics[width=\linewidth]{ui-robot-hao-hung.png}\\[3pt]
{\small\textbf{(b)} Hào hứng (\texttt{excited}): phân tích xong}
\end{minipage}

\vspace{0.8em}
\begin{minipage}[t]{0.485\textwidth}
\centering
\includegraphics[width=\linewidth]{ui-robot-lo-lang.png}\\[3pt]
{\small\textbf{(c)} Lo lắng (\texttt{concerned}): có nguy cơ mức cao}
\end{minipage}\hfill
\begin{minipage}[t]{0.485\textwidth}
\centering
\includegraphics[width=\linewidth]{ui-robot-dong-cam.png}\\[3pt]
{\small\textbf{(d)} Đồng cảm (\texttt{empathetic}): độ tin cậy không xác định}
\end{minipage}
\caption[Robot Vision ở bốn trạng thái biểu cảm]{Ảnh minh hoạ do người dùng cung cấp cho bốn trạng thái robot Vision được hỗ trợ trong source.}
\label{fig:robot-trangthai}
\end{figure}

Ngoài các trang chức năng, giao diện có nhóm công cụ trợ năng và điều khiển bổ
trợ như Bảng~\ref{tab:congcu-trogiup}. Nhóm này giúp người dùng nghe kết quả,
điều khiển bằng giọng nói và quản lý phiên làm việc. Các công cụ chỉ hỗ trợ
tiếp cận thông tin; cảnh báo an toàn vẫn dựa trên ảnh tĩnh và không thay thế
thiết bị hỗ trợ di chuyển đã được kiểm chứng.

\begin{table}[H]
\centering
\small\linespread{1.05}\selectfont
\caption{Công cụ trợ năng và điều khiển bổ trợ trên giao diện.}
\label{tab:congcu-trogiup}
\begin{tabular*}{\textwidth}{@{\extracolsep{\fill}}>{\raggedright\arraybackslash\bfseries\color{daunavy}}p{3.3cm} >{\raggedright\arraybackslash}p{3.0cm} >{\raggedright\arraybackslash}p{8.0cm}}
\toprule
\rowcolor{gray100}
Công cụ & Vị trí & Tác dụng \\
\midrule
Chế độ trợ năng & Thanh bên & Tăng cỡ chữ và nút, làm rõ viền focus khi dùng bàn phím. \\
Bật giọng nói & Thanh bên, Cài đặt & Bật/tắt các nút đọc kết quả thành tiếng. \\
Trò chuyện với Vision & Thanh bên, robot & Mở khung chat Vision để hỏi theo ngữ cảnh trang và ảnh gần nhất, bằng chữ hoặc ghi âm. \\
Nghe kết quả, nghe cảnh báo & Sau khi phân tích & Đọc cảnh báo trước, rồi đến phần mô tả và các nội dung còn lại. \\
Tìm một vật trong ảnh & Sau khi phân tích & Nhập tên vật cần tìm; Vision trả lời theo ảnh hiện có, kết quả không phải định vị chính xác. \\
Lệnh giọng nói & Trợ lý giọng nói & Điều hướng trang (``mở camera'', ``mở OCR''), ``tìm cái cốc'', ``đọc lại kết quả'', ``trợ giúp''; chỉ xử lý sau khi người dùng ghi âm và bấm nhận diện. \\
Xuất báo cáo TXT & Sau khi phân tích & Tải tệp TXT gồm mô tả, vật thể, cảnh báo, OCR, tiền, gợi ý và điều hướng. \\
Kiểm tra kết nối Gemini & Cài đặt & Gửi một yêu cầu nhỏ để xác nhận khoá và quyền truy cập model. \\
Xoá phiên hiện tại & Thanh bên & Xoá kết quả, ảnh và hội thoại đang giữ trong phiên. \\
\bottomrule
\end{tabular*}
\end{table}

\section{Kết quả chạy thử và minh chứng}

Các ảnh dưới đây được chụp từ những lần chạy giao diện thực tế. Với ảnh tiền,
Gemini hiển thị kết quả nhận diện tờ 200.000 đồng, nhưng bản ghi lịch sử của
lượt này vẫn có \texttt{money.detected = false}. Trang Thống kê hiện báo một
lượt nhận diện tiền; đối chiếu lịch sử cho thấy lượt đó đến từ OCR cục bộ đọc
sai văn bản trên tờ tiền, không phải kết quả Gemini. Vì vậy, chỉ số này chưa
đáng tin cậy. Kết quả hỗ trợ di chuyển chỉ dựa trên ảnh tĩnh và không phải
chỉ dẫn an toàn để sử dụng ngoài đời.

\begin{figure}[H]
\centering
\begin{minipage}[t]{0.485\textwidth}
\centering
\includegraphics[width=\linewidth]{screenshot-gemini-nhan-dien-tien.png}\\[3pt]
{\small (a) Lần chạy chế độ nhận diện tiền Việt.}
\end{minipage}\hfill
\begin{minipage}[t]{0.485\textwidth}
\centering
\includegraphics[width=\linewidth]{screenshot-gemini-ho-tro-di-chuyen.png}\\[3pt]
{\small (b) Cảnh báo từ ảnh giao thông tĩnh.}
\end{minipage}
\caption{Ảnh chụp các lần phân tích ảnh bằng Gemini.}
\label{fig:gemini-thucte}
\end{figure}

Ảnh phong cảnh dùng cho hai chế độ dưới đây được lấy từ Wikimedia Commons;
tác giả là Acroterion, giấy phép CC BY-SA 4.0
(\url{https://commons.wikimedia.org/wiki/File:Catoctin_Mountain_and_farm_MD1.jpg}).

\begin{figure}[H]
\centering
\begin{minipage}[t]{0.485\textwidth}
\centering
\includegraphics[width=\linewidth]{screenshot-gemini-mo-ta-toan-canh.png}\\[3pt]
{\small (a) Mô tả toàn cảnh phong cảnh.}
\end{minipage}\hfill
\begin{minipage}[t]{0.485\textwidth}
\centering
\includegraphics[width=\linewidth]{screenshot-gemini-nhan-biet-vat-the.png}\\[3pt]
{\small (b) Nhận biết vật thể và vị trí tương đối.}
\end{minipage}
\caption{Hai chế độ Vision chạy trên ảnh phong cảnh công khai; vị trí là mô tả tương đối trong ảnh, không phải định vị chính xác.}
\label{fig:mota-vatthe}
\end{figure}

Trong thử nghiệm hỏi đáp và tìm vật, mô hình trả về JSON thô ở giao diện.
Đây là vấn đề định dạng đầu ra quan sát được trong lần chạy thử, không phải
lỗi kết nối Gemini. Chức năng TTS đã tạo được trình phát âm thanh; camera và
STT mới được kiểm tra ở mức giao diện, chưa thử bằng thiết bị camera hoặc
microphone thực.

\begin{figure}[H]
\centering
\begin{minipage}[t]{0.485\textwidth}
\centering
\includegraphics[width=\linewidth]{screenshot-gemini-hoi-dap-anh.png}\\[3pt]
{\small (a) Hỏi đáp ảnh; nội dung trả về đang ở dạng JSON.}
\end{minipage}\hfill
\begin{minipage}[t]{0.485\textwidth}
\centering
\includegraphics[width=\linewidth]{screenshot-tim-vat-trong-anh.png}\\[3pt]
{\small (b) Giao diện tìm vật trong ảnh.}
\end{minipage}
\caption{Minh chứng giao diện hỏi đáp ảnh và tìm vật.}
\label{fig:hoidap-timvat}
\end{figure}

\begin{figure}[H]
\centering
\begin{minipage}[t]{0.485\textwidth}
\centering
\includegraphics[width=\linewidth]{screenshot-ocr-result-vie.png}\\[3pt]
{\small (a) OCR cục bộ với dữ liệu tiếng Việt tạm thời; kết quả còn sai.}
\end{minipage}\hfill
\begin{minipage}[t]{0.485\textwidth}
\centering
\includegraphics[width=\linewidth]{screenshot-thong-ke.png}\\[3pt]
{\small (b) Trang thống kê sau các lần chạy thử.}
\end{minipage}
\caption{Kết quả OCR cục bộ và thống kê; chỉ số nhận diện tiền chưa khớp bản ghi phân tích.}
\label{fig:ocr-thongke}
\end{figure}

\begin{figure}[H]
\centering
\includegraphics[width=0.68\textwidth]{screenshot-gemini-ho-tro-khiem-thi.png}
\caption{Chế độ hỗ trợ người khiếm thị đọc nội dung tờ tiền bằng Gemini; phản hồi chỉ là kết quả của một ảnh thử.}
\label{fig:hotro-khiemthi}
\end{figure}

\begin{figure}[H]
\centering
\begin{minipage}[t]{0.485\textwidth}
\centering
\includegraphics[width=\linewidth]{screenshot-tts-giong-doc.png}\\[3pt]
{\small (a) Trình phát TTS được tạo sau thao tác đọc.}
\end{minipage}\hfill
\begin{minipage}[t]{0.485\textwidth}
\centering
\includegraphics[width=\linewidth]{screenshot-lich-su-phan-tich.png}\\[3pt]
{\small (b) Trang lịch sử phân tích.}
\end{minipage}
\caption{Minh chứng TTS và trang lịch sử của ứng dụng.}
\label{fig:tts-lichsu}
\end{figure}

\begin{figure}[H]
\centering
\begin{minipage}[t]{0.485\textwidth}
\centering
\includegraphics[width=\linewidth]{screenshot-giao-dien-camera.png}\\[3pt]
{\small (a) Giao diện camera; chưa xác nhận chụp bằng phần cứng.}
\end{minipage}\hfill
\begin{minipage}[t]{0.485\textwidth}
\centering
\includegraphics[width=\linewidth]{screenshot-giao-dien-tro-ly-giong-noi.png}\\[3pt]
{\small (b) Giao diện trợ lý giọng nói; chưa thử micro thật.}
\end{minipage}
\caption{Hai giao diện cần thiết bị ngoại vi để kiểm thử đầy đủ.}
\label{fig:camera-stt-ui}
\end{figure}

\begin{figure}[H]
\centering
\begin{minipage}[t]{0.485\textwidth}
\centering
\includegraphics[width=\linewidth]{screenshot-cai-dat-visionai.png}\\[3pt]
{\small (a) Tuỳ chọn trợ năng, mức chi tiết và giọng đọc.}
\end{minipage}\hfill
\begin{minipage}[t]{0.485\textwidth}
\centering
\includegraphics[width=\linewidth]{screenshot-huong-dan-visionai.png}\\[3pt]
{\small (b) Trang hướng dẫn sử dụng.}
\end{minipage}
\caption{Minh chứng giao diện cài đặt và hướng dẫn.}
\label{fig:cai-dat-huongdan}
\end{figure}

\section{Kiểm thử hệ thống}

Source có bộ kiểm thử cục bộ \texttt{tests/test\_core.py}. Trong lần kiểm tra
bộ source bàn giao bằng Python 3.13.15, lệnh
\texttt{python -m unittest discover -s tests -v} đạt 15/15 ca; ứng dụng
Streamlit khởi động và endpoint health trả \texttt{ok}. Giao diện đã được
kiểm tra thủ công với một số luồng chức năng; kết quả chạy Gemini, OCR và TTS
được trình bày ở trên. Các kiểm thử unit không gọi Gemini thật, không đo độ
chính xác mô hình và không thay thế đánh giá người dùng.

\begin{table}[H]
\centering
\small
\caption{Mười lăm kiểm thử unit có trong \texttt{tests/test\_core.py}.}
\label{tab:kiemthu}
\begin{tabular*}{\textwidth}{@{\extracolsep{\fill}}>{\bfseries\color{daunavy}}p{0.8cm} p{7.0cm} p{6.0cm}}
\toprule
\rowcolor{gray100}
Mã & Hành vi được kiểm tra & Kết quả mong đợi trong test \\
\midrule
UT01 & Chuẩn hoá ảnh PNG thành data URL JPEG & URL kết quả có MIME type JPEG. \\
UT02 & Từ chối dữ liệu không phải ảnh & Phát sinh \texttt{ValueError}. \\
UT03 & Từ chối tệp ảnh rỗng & Phát sinh \texttt{ValueError}. \\
UT04 & Lưu 105 mục lịch sử & JSON hợp lệ, giữ tối đa 100 mục mới nhất. \\
UT05 & Tạo câu trả lời với ngữ cảnh trang/hội thoại & Prompt mock chứa câu hỏi và ngữ cảnh; hàm trả lời đúng. \\
UT06 & Từ chối câu hỏi Vision rỗng & Phát sinh lỗi trước khi gọi Gemini. \\
UT07 & Gửi ảnh kèm câu hỏi Vision & Request mock có cả phần văn bản và ảnh. \\
UT08 & Tạo lời đọc cho hỗ trợ di chuyển & Có lưu ý ảnh tĩnh và nội dung điều hướng. \\
UT09 & Tạo lời đọc kết quả có cảnh báo/OCR & Bao gồm nguy cơ, văn bản, gợi ý và cảnh báo. \\
UT10 & Tạo lời đọc nhiều dòng OCR & Giữ thứ tự dòng và bỏ qua dòng trống. \\
UT11 & Tập trạng thái robot & Bao gồm các trạng thái trợ năng bổ sung. \\
UT12 & Phân tích lệnh thoại & Nhận điều hướng/tìm vật/đọc lại; câu hỏi thường không thành lệnh. \\
UT13 & Tạo chỉ dẫn mô tả ưu tiên an toàn & Prompt có lưu ý ảnh tĩnh và không khẳng định an toàn. \\
UT14 & Cập nhật yêu thích & Trạng thái lưu lại và khóa định danh mục không đổi. \\
UT15 & Chọn giọng Edge TTS & Dùng giọng được yêu cầu, không đổi giọng mặc định. \\
\bottomrule
\end{tabular*}
\end{table}

Bảng~\ref{tab:kiemthu} ghi lại phạm vi các unit test đã chạy. Những test gửi
yêu cầu Gemini dùng client giả lập; kiểm thử này không xác nhận API Gemini,
Tesseract, thiết bị camera/micro hay dịch vụ TTS thực tế.

\section{Đánh giá bước đầu}

Phiên bản source đã có chuỗi chức năng chính từ nhận ảnh/âm thanh, phân tích,
chuẩn hoá kết quả đến phản hồi và lưu lịch sử. Kết quả 15/15 chỉ xác nhận các
kiểm thử phần mềm cục bộ đã viết; việc đánh giá định lượng độ chính xác OCR,
Vision, STT, độ trễ và mức độ phù hợp với người dùng thực tế cần bộ dữ liệu,
kịch bản riêng và kiểm thử có sự tham gia của người dùng.

%% ---------------------------------------------------------------------
%  CHƯƠNG 5 — KẾT QUẢ VÀ ĐÁNH GIÁ
% ---------------------------------------------------------------------
\chapter{Kết quả và đánh giá}
\label{ch:ketqua}

\begin{muctieu}
Chương này tổng hợp chức năng của source VisionAI Pro và kết quả kiểm tra
đã thực hiện. Kiểm thử phần mềm được tách biệt với đánh giá độ chính xác,
độ trễ và mức độ an toàn của mô hình AI.
\end{muctieu}

\section{Kết quả đạt được}

Qua đối chiếu mã nguồn và chạy thử ứng dụng, các chức năng hiện có gồm:
\begin{itemize}[leftmargin=1.6em, itemsep=2pt]
  \item Giao diện Streamlit có 11 trang, robot Vision đồ hoạ tương tác và lối tắt tới các công cụ chính.
  \item Phân tích ảnh bằng Gemini theo danh sách model cấu hình được; hỗ trợ mô tả cảnh, vật thể, nguy cơ, hỗ trợ di chuyển và nhận diện tiền. Khi Gemini không khả dụng, source thử OCR cục bộ và cảnh báo rằng không phân tích được cảnh/vật thể.
  \item OCR tiếng Việt/Anh bằng Tesseract cục bộ; không cần gọi Gemini cho chế độ OCR.
  \item Chuẩn hoá ảnh đầu vào: giới hạn 15 MB, kiểm tra ảnh hợp lệ, sửa hướng EXIF, đổi RGB, thu nhỏ tối đa 1600 pixel và nén JPEG.
  \item Escape nội dung do AI tạo trước khi chèn vào HTML và hiển thị lỗi ảnh/dịch vụ bằng thông báo dễ hiểu.
  \item Hỏi đáp theo ảnh gần nhất, hỏi đáp trong khung chat Vision theo ngữ cảnh trang, tìm một vật trong ảnh và lệnh thoại điều hướng/đọc lại.
  \item Nhận dạng giọng nói qua SpeechRecognition/Google Web Speech sau khi người dùng ghi âm và xác nhận; Edge TTS và Windows SAPI5/pyttsx3 đọc kết quả.
  \item Chế độ trợ năng, lựa chọn giọng đọc, mức chi tiết, kiểm tra kết nối Gemini, lịch sử có lọc/yêu thích/xoá xác nhận, thống kê và xuất TXT.
  \item Lịch sử cục bộ được giới hạn 100 bản ghi và ghi theo cách thay thế tệp nguyên tử.
\end{itemize}

\section{Kết quả kiểm thử thực tế}

\begin{table}[H]
\centering
\caption{Kết quả kiểm tra phiên bản source trong môi trường hiện có.}
\label{tab:ketqua}
\begin{tabular*}{\textwidth}{@{\extracolsep{\fill}}>{\bfseries\color{daunavy}}p{3.1cm} p{5.5cm} p{5.1cm}}
\toprule
\rowcolor{gray100}
Hạng mục & Cách kiểm tra & Kết quả \\
\midrule
Kiểm thử đơn vị & Python 3.13.15; chạy \texttt{python -m unittest discover -s tests -v} trên source bàn giao & 15/15 kiểm thử đạt. \\
Khởi chạy ứng dụng & Chạy Streamlit từ source bàn giao & Máy chủ khởi động; endpoint kiểm tra sức khoẻ trả về \texttt{ok}. \\
Giao diện & Chạy thủ công một số luồng trên trình duyệt, gồm ảnh, OCR, hỏi đáp, lịch sử và thống kê & Các trang đã mở được; camera và STT chưa thử bằng thiết bị thật. \\
Gemini API thật & Kiểm tra kết nối và phân tích ảnh tiền, ảnh giao thông bằng \texttt{gemini-3.5-flash-lite} & Kết nối thành công; ảnh tiền được mô tả là 200.000 đồng (0,98), ảnh giao thông có cảnh báo nguy cơ cao (0,90). \\
Hỏi đáp/tìm vật & Gửi câu hỏi theo ảnh và yêu cầu tìm vật & Gemini phản hồi; giao diện hiển thị JSON thô thay vì câu trả lời tự nhiên. \\
OCR cục bộ & Chạy Tesseract trên ảnh tờ tiền mẫu & Có kết quả nhưng nhiều chữ tiếng Việt sai; máy thiếu dữ liệu \texttt{vie.traineddata}. \\
TTS & Yêu cầu đọc kết quả trong ứng dụng & Tạo được trình phát âm thanh; chưa đánh giá chất lượng giọng trên thiết bị người dùng. \\
Camera/STT & Mở các trang camera và trợ lý giọng nói & Chưa kiểm tra chụp ảnh bằng camera thật hoặc ghi âm bằng microphone thật. \\
Độ chính xác/độ trễ & Chưa có bộ dữ liệu gán nhãn và phép đo chuyên biệt & Chưa có số liệu; không suy luận từ số ca unit test. \\
\bottomrule
\end{tabular*}
\end{table}

Bảng~\ref{tab:ketqua} phân biệt unit test với các lần chạy dịch vụ thật; không
diễn giải 15/15 unit test thành độ chính xác của mô hình AI. Các unit test
bao phủ chuẩn hoá và từ chối ảnh lỗi, lưu lịch sử/đánh dấu yêu thích, tạo nội
dung đọc thành tiếng, trạng thái robot, lệnh thoại và cách tạo yêu cầu Gemini.
Các lần gọi Gemini xác nhận được một số luồng trên hai ảnh mẫu, nhưng không
đủ để khái quát chất lượng. OCR cục bộ cho kết quả tiếng Việt kém trong cấu
hình máy hiện tại vì thiếu dữ liệu \texttt{vie}; TTS tạo được trình phát.
Camera và STT chưa được thử bằng phần cứng thật. Bản ghi nhận diện tờ tiền
có mô tả 200.000 đồng nhưng trường \texttt{money.detected} là false, nên
chỉ số nhận diện tiền trong trang Thống kê không tăng.

\section{Bàn luận}

Ứng dụng kết hợp dịch vụ AI trực tuyến với xử lý cục bộ có phạm vi rõ ràng:
Gemini thực hiện phân tích thị giác và hỏi đáp; Tesseract đảm nhiệm OCR khi
được chọn; Google Web Speech nhận dạng âm thanh; Edge TTS tạo giọng đọc và
Windows SAPI5 là phương án dự phòng. Vì vậy, không nên mô tả STT hiện tại là
chạy cục bộ. Chế độ trợ năng, robot Vision và lệnh thoại tạo thêm cách tương
tác, nhưng người dùng vẫn cần điều khiển giao diện và chủ động xác nhận thu âm.

Các biện pháp giới hạn kích thước, chuẩn hoá ảnh, ghi lịch sử nguyên tử,
thử model theo danh sách và xử lý lỗi giúp tăng tính bền vững của nguyên mẫu.
Tuy vậy, các biện pháp này không chứng minh độ chính xác của AI. Đặc biệt,
kết quả không phát hiện nguy cơ không đồng nghĩa môi trường an toàn; hỗ trợ
di chuyển chỉ dựa trên một ảnh tĩnh.

\section{Hạn chế}

\begin{itemize}[leftmargin=1.6em, itemsep=2pt]
  \item Gemini cần API key, kết nối mạng và hạn mức phù hợp; danh sách model dự phòng không khắc phục được khoá sai hoặc mất mạng.
  \item STT sử dụng Google Web Speech và cần Internet; Edge TTS cũng cần mạng. Windows SAPI5 chỉ là phương án TTS dự phòng.
  \item Tesseract và dữ liệu ngôn ngữ \texttt{vie}/\texttt{eng} cần được cài đặt riêng.
  \item Camera chỉ chụp ảnh theo thao tác người dùng; hỗ trợ di chuyển không theo dõi video, không đo khoảng cách và không thay thế gậy, chó dẫn đường hay thiết bị đã được kiểm chứng.
  \item Robot là đồ hoạ CSS/SVG tạo cảm giác chuyển động, không phải mô hình 3D có thể xoay; hoạt ảnh không đồng bộ khẩu hình với tiếng nói.
  \item Giọng nói theo luồng ghi âm--chuyển văn bản--gửi Gemini--đọc phản hồi, chưa phải hội thoại âm thanh thời gian thực.
  \item Cá nhân hoá giới hạn ở trạng thái phiên; không có tài khoản, hồ sơ người dùng lâu dài, cơ sở dữ liệu hoặc cấu hình production nhiều người dùng.
  \item Lịch sử lưu cục bộ tối đa 100 bản ghi; chưa có bộ dữ liệu để đo độ chính xác, độ trễ hoặc đánh giá có người dùng.
\end{itemize}

\section{Kết luận đánh giá}

Trên source bàn giao, ứng dụng khởi chạy, health check đạt và bộ kiểm thử
cục bộ đạt 15/15; kết nối Gemini, hai lần phân tích ảnh và TTS cũng đã chạy
được trong các thử nghiệm nêu trên. Những lần chạy này không xác nhận chất
lượng nhận diện tổng quát hay tính an toàn khi sử dụng. Cần cài đủ dữ liệu
OCR tiếng Việt, kiểm thử camera/STT trên thiết bị phù hợp và xây dựng bộ dữ
liệu có nhãn trước khi đưa ra kết luận định lượng.

%% ---------------------------------------------------------------------
%  CHƯƠNG 6 — KẾT LUẬN VÀ HƯỚNG PHÁT TRIỂN
% ---------------------------------------------------------------------
\chapter{Kết luận và hướng phát triển}
\label{ch:ketluan}

\begin{muctieu}
Chương này tổng kết hệ thống đã xây dựng dựa trên source VisionAI Pro,
rút ra bài học và đề xuất các hướng phát triển phù hợp với kiến trúc hiện tại.
\end{muctieu}

\section{Kết luận}

Đồ án đã xây dựng một ứng dụng trợ lý thị giác AI bằng Python và Streamlit,
cho phép tải/chụp ảnh, phân tích bằng Gemini Vision và đọc văn bản cục bộ bằng
Tesseract. Ứng dụng bổ sung hỏi đáp theo ảnh, tìm vật, hỗ trợ di chuyển từ ảnh
tĩnh, lệnh thoại và phản hồi bằng giọng nói. SpeechRecognition/Google Web Speech
được dùng cho STT; Edge TTS và Windows SAPI5 đảm nhiệm TTS. Robot Vision, chế độ
trợ năng, lịch sử có lọc/yêu thích, thống kê, xuất TXT và kiểm tra kết nối
Gemini hoàn thiện các luồng sử dụng chính.

Kiến trúc module tách Vision, OCR, Speech, trợ năng, lệnh thoại và robot.
Trên source bàn giao, ứng dụng khởi chạy, health check đạt và 15/15 kiểm thử
cục bộ đạt. Đây là kết quả kiểm thử phần mềm, không phải số đo độ chính xác
của AI; Gemini API thật và chất lượng sử dụng với người khiếm thị chưa được
đánh giá. Cá nhân hoá hiện chỉ ở mức cấu hình và trạng thái phiên làm việc.

\section{Bài học rút ra}

Quá trình xây dựng cho thấy một ứng dụng AI hỗ trợ người khiếm thị không chỉ
cần mô hình AI mà còn cần thiết kế luồng dữ liệu, xử lý lỗi, khả năng tiếp cận,
phản hồi âm thanh và chỉ dẫn an toàn. Việc tách OCR cục bộ, nhận dạng giọng nói
trực tuyến, TTS, trợ năng và giao diện robot thành các module riêng giúp dễ
kiểm thử và thay thế công nghệ.

Đồ án cũng củng cố kỹ năng Python, Streamlit, xử lý ảnh, tích hợp API AI,
xử lý âm thanh, quản lý biến môi trường và thiết kế một pipeline đa phương thức.

\section{Hướng phát triển}

Một số hướng phát triển phù hợp với source hiện tại gồm:
\begin{itemize}[leftmargin=1.6em, itemsep=2pt]
  \item Xây dựng hồ sơ người dùng và cơ chế cá nhân hoá lâu dài thay cho cấu hình chỉ tồn tại trong phiên.
  \item Bổ sung bộ dữ liệu kiểm thử có nhãn để đo định lượng OCR, Vision, STT và khả năng nhận diện tiền.
  \item Nghiên cứu STT cục bộ hoặc các lựa chọn nhận dạng khác để giảm phụ thuộc vào dịch vụ mạng, đồng thời đánh giá quyền riêng tư.
  \item Mở rộng khả năng xử lý ảnh/video thời gian thực cho cảnh báo vật cản.
  \item Bổ sung cơ chế cache, theo dõi lỗi mạng và phương án vận hành phù hợp khi Gemini, Google Web Speech hoặc TTS trực tuyến không khả dụng.
  \item Đánh giá khả năng tiếp cận, lệnh thoại, kích thước chữ/nút và nội dung cảnh báo với người dùng mục tiêu.
  \item Phát triển phiên bản mobile hoặc thiết bị đeo sau khi hoàn thiện pipeline và kiểm thử với người dùng thực tế.
\end{itemize}


% ---------------------------------------------------------------------
%  TÀI LIỆU THAM KHẢO VÀ PHỤ LỤC
% ---------------------------------------------------------------------
\backmatter

\cleardoublepage
\phantomsection
\addcontentsline{toc}{chapter}{Tài liệu tham khảo}
\bibliographystyle{plain}
\bibliography{refs}

% ---------------------------------------------------------------------
%  PHỤ LỤC
% ---------------------------------------------------------------------
\chapter{Phụ lục}
\markboth{Phụ lục}{}

\section*{A. Hướng dẫn cài đặt và chạy chương trình}
\phantomsection
\addcontentsline{toc}{section}{A. Hướng dẫn cài đặt và chạy chương trình}

Các bước dưới đây được lấy theo quy trình trong \texttt{README.md} của source
VisionAI Pro. Mở PowerShell tại thư mục chứa \texttt{app.py}.

\begin{lstlisting}[language=bash, caption={Tạo môi trường và chạy ứng dụng trên Windows.}]
py -3.11 -m venv .venv
.\.venv\Scripts\Activate.ps1
python -m pip install -r requirements.txt
Copy-Item .env.example .env
# Mo .env va tu dien GEMINI_API_KEY tren may cua ban.
python -m streamlit run app.py
\end{lstlisting}

Có thể chạy \texttt{CHAY\_VISIONAI.bat} thay cho lệnh Streamlit. OCR cục bộ
cần cài Tesseract và dữ liệu ngôn ngữ \texttt{vie}/\texttt{eng}; khi cần có
thể dùng \texttt{setup\_ocr.ps1}. Chạy kiểm thử phần mềm bằng:

\begin{lstlisting}[language=bash, caption={Chạy các kiểm thử cục bộ không cần gọi Gemini.}]
.\.venv\Scripts\python.exe -m unittest discover -s tests -v
\end{lstlisting}

\subsection*{Cấu hình \texttt{.env}}

Có thể tham khảo các biến sau trong \texttt{.env.example}:

\begin{lstlisting}[language=bash, caption={Các biến môi trường của hệ thống.}]
GEMINI_API_KEY=PASTE_YOUR_GEMINI_API_KEY_HERE
GEMINI_MODELS=gemini-3.8-flash,gemini-3.6-flash,gemini-3.5-flash-lite
EDGE_TTS_VOICE=vi-VN-HoaiMyNeural
EDGE_TTS_RATE=+0%
EDGE_TTS_PITCH=+0Hz
STT_LANGUAGE=vi-VN
TESSERACT_CMD=
\end{lstlisting}

Đây là tên biến và giá trị mặc định mẫu; không đưa API key thật vào báo cáo,
ảnh chụp hoặc kho mã nguồn.

\section*{B. Cấu trúc source tham khảo}
\phantomsection
\addcontentsline{toc}{section}{B. Cấu trúc source tham khảo}

\begin{lstlisting}[caption={Cấu trúc chính của source VisionAI Pro.}]
tro-ly-thi-giac/
|-- app.py
|-- requirements.txt
|-- .env.example
|-- README.md
|-- setup_ocr.ps1
|-- check_env.py
|-- data/
|   `-- history.json
|-- tests/
|   `-- test_core.py
|-- ui/
|   `-- vision_robot/
`-- modules/
    |-- ai_vision.py
    |-- local_ocr.py
    |-- accessibility.py
    |-- robot_ui.py
    |-- voice_commands.py
    |-- speech.py
    `-- utils.py
\end{lstlisting}

\section*{C. Bảng tự kiểm trước khi nộp}
\phantomsection
\addcontentsline{toc}{section}{C. Bảng tự kiểm trước khi nộp}

\begin{tukiem}
  \item Kiểm tra \texttt{.env} không chứa API key thật khi nộp hoặc chia sẻ source công khai.
  \item Kiểm tra Tesseract và dữ liệu \texttt{vie}/\texttt{eng} nếu trình diễn OCR tiếng Việt.
  \item Kiểm tra \texttt{pip install -r requirements.txt} và \texttt{python -m streamlit run app.py} trên máy trình diễn.
  \item Kiểm tra robot, chế độ trợ năng, tải ảnh, camera, OCR, Vision, hỏi đáp, lệnh thoại, STT, TTS, lịch sử, yêu thích và xuất TXT.
  \item Xác nhận STT qua Google Web Speech và Edge TTS cần Internet; không mô tả STT hiện tại là Faster-Whisper cục bộ.
  \item Kiểm tra xử lý ảnh lớn/hỏng, lệnh xoá lịch sử có xác nhận và các lỗi cấu hình Gemini.
  \item Mọi số liệu đánh giá trong báo cáo phải lấy từ kết quả kiểm thử thực tế.
  \item Kiểm thử Gemini thật, độ chính xác mô hình và trải nghiệm người dùng cần được ghi riêng; 15 kiểm thử unit không thay thế các đánh giá này.
  \item Không mô tả Faster-Whisper, database hoặc xử lý video thời gian thực như chức năng đã triển khai.
\end{tukiem}

\section*{D. Nhận xét của giảng viên hướng dẫn}
\phantomsection
\addcontentsline{toc}{section}{D. Nhận xét của giảng viên hướng dẫn}

\khungnhanxet[7.5cm]

\vspace{0.8em}
\noindent\begin{minipage}[t]{0.5\textwidth}
\color{daunavy}
{\color{daumaroon}\bfseries Điểm số}

\vspace{0.7em}
\dongchamnhan{Bằng số:}

\vspace{0.9em}
\dongchamnhan{Bằng chữ:}
\end{minipage}\hfill
\okyten{\dauGVHD}{Giảng viên hướng dẫn}


\end{document}